\pdfoutput=1
\RequirePackage{ifpdf}
\ifpdf % We~are running pdfTeX in pdf mode
\documentclass[pdftex]{sigma}
\else
\documentclass{sigma}
\fi

\usepackage{upgreek}

\numberwithin{equation}{section}

\newtheorem{Theorem}{Theorem}[section]
\newtheorem*{Theorem*}{Theorem}
\newtheorem{Corollary}[Theorem]{Corollary}
\newtheorem{Lemma}[Theorem]{Lemma}
\newtheorem{Proposition}[Theorem]{Proposition}
 { \theoremstyle{definition}
\newtheorem{Definition}[Theorem]{Definition}
\newtheorem{Example}[Theorem]{Example}
\newtheorem{Remark}[Theorem]{Remark} }

\begin{document}

\allowdisplaybreaks

\newcommand{\arXivNumber}{2509.17976}

\renewcommand{\PaperNumber}{089}

\FirstPageHeading

\ShortArticleName{Integrable Twelve-Component Nonlinear Dynamical System}

\ArticleName{Integrable Twelve-Component Nonlinear Dynamical \\ System on a Quasi-One-Dimensional Lattice}

\Author{Oleksiy O.~VAKHNENKO~$^{\rm a}$ and Vyacheslav O.~VAKHNENKO~$^{\rm b}$}

\AuthorNameForHeading{O.O.~Vakhnenko and V.O.~Vakhnenko}

\Address{$^{\rm a)}$~Bogolyubov Institute for Theoretical Physics of NAS of Ukraine,\\
\hphantom{$^{\rm a)}$}~14-B Metrologichna Str., Ky\"{\i}v 03143, Ukraine}
\EmailD{\href{mailto:vakhnenko@bitp.kiev.ua}{vakhnenko@bitp.kiev.ua}}

\Address{$^{\rm b)}$~Subbotin Institute of Geophysics of NAS of Ukraine,\\
\hphantom{$^{\rm b)}$}~63-B Bohdan Khmel'nyts'kyy Str., Ky\"{\i}v 01054, Ukraine}
\EmailD{\href{mailto:vakhnenko@ukr.net}{vakhnenko@ukr.net}}

\ArticleDates{Received May 24, 2025, in final form September 17, 2025; Published online October 22, 2025}

\Abstract{Bearing in mind the potential physical applicability of multicomponent completely integrable nonlinear dynamical models on quasi-one-dimensional lattices we have developed the novel twelve-component and six-component semi-discrete nonlinear inregrable systems in the framework of semi-discrete Ablowitz--Kaup--Newell--Segur scheme. The set of lowest local conservation laws found by the generalized direct recurrent technique was shown to be indispensable constructive tool in the reduction procedure from the prototype to actual field variables. Two types of admissible symmetries for the twelve-component system and one type of symmetry for the six-component system have been established. The mathematical structure of total local current was shown to support the charge transportation only by four of six subsystems incorporated into the twelve-component system under study. The twelve-component system is able to model the actions of external parametric drive and external uniform magnetic field via time dependencies and phase factors of coupling parameters.\looseness=-1}

\Keywords{Lax integrability; quasi-one-dimensional lattice; multicomponent system; nonlinear dynamics; $\mathcal{PT}$ symmetry}

\Classification{39A36; 37K10; 35Q55; 58J70}

\section{Introduction}
%\label{sec1}

Since the fundamental works by Toda \cite{toda-JPSJ-22-431, toda-JPSJ-23-501, toda-PhR-18-1} as well as by Ablowitz and Ladik \mbox{\cite{ ablowitz-SAM-58-17,ablowitz-SAM-55-213, ablowitz-JMP-17-1011}} the integrable nonlinear dynamical systems on one-dimensional and quasi-one-dimensional lattices
are steadily gaining a significant influence on modeling a wide variety of diverse nonlinear phenomena in physical, biological and applied sciences. In this respect, the multicomponent differential-difference (semi-discrete) nonlinear integrable systems \cite{ablowitz-PLA-253-287, ablowitz-CSR-11-159,gerdzhikov-TMP-52-676, tsuchida-JMP-39-4785, tsuchida-JPA-32-2239, vakhnenko-JPA-32-5735, vakhnenko-PRE-60-R2492,vakhnenko-JNMP-24-250, vakhnenko-PRE-61-7110, vakhnenko-EPJP-137-1176} acquire an important place due to their physical importance and mathematical reliability.

Sometimes, however, the lack of physical imagination inspires the semi-discrete nonlinear integrable systems allegedly claimed to be multicomponent but actually composed of several physically uncoupled one-component or two-component basic systems \cite{babalic-RJP-63-114, babalic-IJMPB-34-2050274,babalic-JPA-50-415201,carstea-JPA-48-055205, wen-AML-123-107591,yu-WM-94-102500, yuan-CSF-168-113180}. The critical overview of semi-discrete nonlinear integrable systems characterized by the false multicomponentness is given in our recent paper \cite{vakhnenko-PLA-452-128460}.

The progress in the development of new differential-differential multicomponent nonlinear integrable systems has been prompted by the guiding rules known as the Ablowitz--Kaup--Newell--Segur scheme \cite{ablowitz-PRL-31-125, ablowitz-SAM-53-249}. The spatially discretized version of Ablowitz--Kaup--Newell--Segur rules is now widely recognized as the prospective tool for the construction of new differential-difference (semi-discrete) multicomponent nonlinear integrable systems. The key points of spatially discretized Ablowitz--Kaup--Newell--Segur scheme are summarized in our previous papers \cite{vakhnenko-UJP-69-168, vakhnenko-UJP-69-577}. We followed this constructive procedure while developing the integrable twelve-component nonlinear dynamical system on a quasi-one-dimensional lattice suggested in the present paper.

\section[Semi-discrete zero-curvature equation and the mutually consistent ans\"{a}tze for the auxiliary spectral and evolutionary matrices]{Semi-discrete zero-curvature equation and the mutually\\ consistent ans\"{a}tze for the auxiliary spectral\\ and evolutionary matrices}
%\label{sec2}


Any classic nonlinear evolutionary system on an infinite quasi-one-dimensional lattice is said to be integrable in the Lax sense provided it admits a matrix-valued semi-discrete zero-curvature representation \cite{faddeev-HMTS-2007, tu-JPA-23-3903}
\begin{equation}\label{eq-2.1}
	\frac{\mathrm{d}}{\mathrm{d}\tau}L(n|z)=A(n+1|z)L(n|z)-L(n|z)A(n|z)
\end{equation}
with the auxiliary square matrices $L(n|z)$ and $A(n|z)$ referred to as the spectral and evolutionary matrices, respectively.
Here $\tau$ denotes the continuous time variable, $n$ stands for the discrete space variable running through the integer numbers from minus infinity to plus infinity, while~$z$ marks the time- and space-independent spectral parameter. The spectral matrix $L(n|z)$ is assumed to be the non-degenerate one ($\det L(n|z)\neq0$).

In this paper, we suggest the auxiliary matrices $L(n|z)$ and $A(n|z)$ in the following forms:
\begin{gather}\label{eq-2.2}
	L(n|z) = \begin{pmatrix}
		0 & t_{12}(n) & u_{13}(n)z^{-1} & 0 \\
		t_{21}(n) & r_{22}(n)z^{2} + t_{22}(n) & s_{23}(n)z + u_{23}(n)z^{-1} & s_{24}(n)z \\
		u_{31}(n)z^{-1} & s_{32}(n)z + u_{32}(n)z^{-1} & t_{33}(n) + v_{33}(n)z^{-2} & t_{34}(n) \\
		0 & s_{42}(n)z & t_{43}(n) & 0	
	\end{pmatrix},
\\ \label{eq-2.3}
	A(n|z) = \begin{pmatrix}
		c_{11}(n) & c_{12}(n) & d_{13}(n)z^{-1} & 0 \\
		c_{21}(n) & a_{22}(n)z^{2} + c_{22}(n) & b_{23}(n)z + d_{23}(n)z^{-1} & b_{24}(n)z \\
		d_{31}(n)z^{-1} & b_{32}(n)z + d_{32}(n)z^{-1} & c_{33}(n) + e_{33}(n)z^{-2} & c_{34}(n) \\
		0 & b_{42}(n)z & c_{43}(n) & c_{44}(n)
	\end{pmatrix},
\end{gather}
where the ansatz (\ref{eq-2.3}) for the evolutionary matrix $A(n|z)$ is extended by the additional terms $c_{11}(n)$ and $c_{44}(n)$ as compared with the ansatz considered previously \cite{vakhnenko-JNMP-18-401, vakhnenko-JNMP-18-415}.

The adopted ans\"{a}tze (\ref{eq-2.2}) and (\ref{eq-2.3}) in combination with the zero-curvature equation (\ref{eq-2.1}) allow to fix the majority of matrix elements $A_{jk}(n|z)$ of evolutionary matrix $A(n|z)$ in terms of prototype field functions $t_{12}(n)$, $u_{13}(n)$, $t_{21}(n)$, $r_{22}(n)$, $t_{22}(n)$, $s_{23}(n)$, $u_{23}(n)$, $s_{24}(n)$, $u_{31}(n)$, $s_{32}(n)$, $u_{32}(n)$, $t_{33}(n)$, $v_{33}(n)$, $t_{34}(n)$, $s_{42}(n)$, $t_{43}(n)$ as well as to recover the set of primary semi-discrete nonlinear equations for the prototype field functions. Under certain plausible assumptions prompted by the on-site local conservation laws, the primary semi-discrete nonlinear equations are convertible into the one or another semi-discrete nonlinear integrable system of preferable interest. In this sense, the suggested ans\"{a}tze (\ref{eq-2.2}) and (\ref{eq-2.3}) for the auxiliary matrices are proved to be mutually consistent.

Let us note that the choice of proper ansatz for the evolutionary matrix $A(n|z)$ is not unique and it manifests a particular integrable system in an infinite hierarchy related to the generic spectral matrix (\ref{eq-2.2}).


\section{Primary semi-discrete nonlinear equations}
%\label{sec3}

To proceed with developing the semi-discrete nonlinear integrable systems of our present interest, we have to
rely upon the primary (prototype) set of semi-discrete equations
\begin{gather}
	\frac{\mathrm{d}}{\mathrm{d}\tau}t_{12}(n)=c_{11}(n+1)t_{12}(n)+c_{12}(n+1)t_{22}(n)+d_{13}(n+1)s_{32}(n)\nonumber\\
\hphantom{\frac{\mathrm{d}}{\mathrm{d}\tau}t_{12}(n)=}{}
	-t_{12}(n)c_{22}(n)-u_{13}(n)b_{32}(n),\label{eq-3.1}
\\
	\frac{\mathrm{d}}{\mathrm{d}\tau}u_{13}(n)=c_{11}(n+1)u_{13}(n)+c_{12}(n+1)u_{23}(n)+d_{13}(n+1)t_{33}(n)\nonumber\\
\hphantom{\frac{\mathrm{d}}{\mathrm{d}\tau}u_{13}(n)=}{}
	-t_{12}(n)d_{23}(n)-u_{13}(n)c_{33}(n),\label{eq-3.2}
\\
	\frac{\mathrm{d}}{\mathrm{d}\tau}t_{21}(n)=c_{22}(n+1)t_{21}(n)+b_{23}(n+1)u_{31}(n)\nonumber\\
\hphantom{\frac{\mathrm{d}}{\mathrm{d}\tau}t_{21}(n)=}{}
	-t_{21}(n)c_{11}(n)-t_{22}(n)c_{21}(n)-s_{23}(n)d_{31}(n),\label{eq-3.3}
\\
	\frac{\mathrm{d}}{\mathrm{d}\tau}r_{22}(n)=c_{22}(n+1)r_{22}(n)+b_{23}(n+1)s_{32}(n)+b_{24}(n+1)s_{42}(n) \nonumber\\
\hphantom{\frac{\mathrm{d}}{\mathrm{d}\tau}r_{22}(n)=}{}
-r_{22}(n)c_{22}(n)-s_{23}(n)b_{32}(n)-s_{24}(n)b_{42}(n),\label{eq-3.4}
\\
	\frac{\mathrm{d}}{\mathrm{d}\tau}t_{22}(n)=c_{21}(n+1)t_{12}(n)+c_{22}(n+1)t_{22}(n)-t_{21}(n)c_{12}(n)-t_{22}(n)c_{22}(n) \nonumber\\
\hphantom{\frac{\mathrm{d}}{\mathrm{d}\tau}t_{22}(n)=}{}
	+b_{23}(n+1)u_{32}(n)+d_{23}(n+1)s_{32}(n)-s_{23}(n)d_{32}(n)-u_{23}(n)b_{32}(n),\label{eq-3.5}
\\
	\frac{\mathrm{d}}{\mathrm{d}\tau}s_{23}(n)=a_{22}(n+1)u_{23}(n)+c_{22}(n+1)s_{23}(n)-r_{22}(n)d_{23}(n)-t_{22}(n)b_{23}(n) \nonumber\\
\hphantom{\frac{\mathrm{d}}{\mathrm{d}\tau}s_{23}(n)=}{}
	+b_{23}(n+1)t_{33}(n)+b_{24}(n+1)t_{43}(n)-s_{23}(n)c_{33}(n)-s_{24}(n)c_{43}(n),\label{eq-3.6}
\\
	\frac{\mathrm{d}}{\mathrm{d}\tau}u_{23}(n)=c_{21}(n+1)u_{13}(n)+c_{22}(n+1)u_{23}(n)-t_{21}(n)d_{13}(n)-t_{22}(n)d_{23}(n) \nonumber\\
\hphantom{\frac{\mathrm{d}}{\mathrm{d}\tau}u_{23}(n)=}{}
	+b_{23}(n+1)v_{33}(n)+d_{23}(n+1)t_{33}(n)-s_{23}(n)e_{33}(n)-u_{23}(n)c_{33}(n),\label{eq-3.7}
\\
	\frac{\mathrm{d}}{\mathrm{d}\tau}s_{24}(n)=c_{22}(n+1)s_{24}(n)+b_{23}(n+1)t_{34}(n)\nonumber\\
\hphantom{\frac{\mathrm{d}}{\mathrm{d}\tau}s_{24}(n)=}{}
	-t_{22}(n)b_{24}(n)-s_{23}(n)c_{34}(n)-s_{24}(n)c_{44}(n),\label{eq-3.8}
\\
	\frac{\mathrm{d}}{\mathrm{d}\tau}u_{31}(n)=d_{32}(n+1)t_{21}(n)+c_{33}(n+1)u_{31}(n)\nonumber\\
\hphantom{\frac{\mathrm{d}}{\mathrm{d}\tau}u_{31}(n)=}{}
	-u_{31}(n)c_{11}(n)-u_{32}(n)c_{21}(n)-t_{33}(n)d_{31}(n), \label{eq-3.9}
\\
	\frac{\mathrm{d}}{\mathrm{d}\tau}s_{32}(n)=b_{32}(n+1)t_{22}(n)+d_{32}(n+1)r_{22}(n)-s_{32}(n)c_{22}(n)-u_{32}(n)a_{22}(n) \nonumber\\
\hphantom{\frac{\mathrm{d}}{\mathrm{d}\tau}s_{32}(n)=}{}
	+c_{33}(n+1)s_{32}(n)+c_{34}(n+1)s_{42}(n)-t_{33}(n)b_{32}(n)-t_{34}(n)b_{42 }(n), \label{eq-3.10}
\\
	\frac{\mathrm{d}}{\mathrm{d}\tau}u_{32}(n)=d_{31}(n+1)t_{12}(n)+d_{32}(n+1)t_{22}(n)-u_{31}(n)c_{12}(n)-u_{32}(n)c_{22}(n) \nonumber\\
\hphantom{\frac{\mathrm{d}}{\mathrm{d}\tau}u_{32}(n)=}{}
	+c_{33}(n+1)u_{32}(n)+e_{33}(n+1)s_{32}(n)-t_{33}(n)d_{32}(n)-v_{33}(n)b_{32}(n),\label{eq-3.11}
\\
	\frac{\mathrm{d}}{\mathrm{d}\tau}t_{33}(n)=b_{32}(n+1)u_{23}(n)+d_{32}(n+1)s_{23}(n)-s_{32}(n)d_{23}(n)-u_{32}(n)b_{23}(n) \nonumber\\
\hphantom{\frac{\mathrm{d}}{\mathrm{d}\tau}t_{33}(n)=}{}
	+c_{33}(n+1)t_{33}(n)+c_{34}(n+1)t_{43}(n)-t_{33}(n)c_{33}(n)-t_{34}(n)c_{43}(n), \label{eq-3.12}
\\
	\frac{\mathrm{d}}{\mathrm{d}\tau}v_{33}(n)=d_{31}(n+1)u_{13}(n)+d_{32}(n+1)u_{23}(n)+c_{33}(n+1)v_{33}(n) \nonumber\\
\hphantom{\frac{\mathrm{d}}{\mathrm{d}\tau}v_{33}(n)=}{}
	-u_{31}(n1)d_{13}(n)-u_{32}(n)d_{23}(n)-v_{33}(n)c_{33}(n), \label{eq-3.13}
\\
	\frac{\mathrm{d}}{\mathrm{d}\tau}t_{34}(n)=d_{32}(n+1)s_{24}(n)+c_{33}(n+1)t_{34}(n)\nonumber\\
\hphantom{\frac{\mathrm{d}}{\mathrm{d}\tau}t_{34}(n)=}{}
	-u_{32}(n)b_{24}(n)-t_{33}(n)c_{34}(n)-t_{34}(n)c_{44}(n), \label{eq-3.14}
\\
	\frac{\mathrm{d}}{\mathrm{d}\tau}s_{42}(n)=b_{42}(n+1)t_{22}(n)+c_{43}(n+1)s_{32}(n)+c_{44}(n+1)s_{42}(n)\nonumber\\
\hphantom{\frac{\mathrm{d}}{\mathrm{d}\tau}s_{42}(n)=}{}
	-s_{42}(n)c_{22}(n)-t_{43}(n)b_{32}(n), \label{eq-3.15}
\\
	\frac{\mathrm{d}}{\mathrm{d}\tau}t_{43}(n)=b_{42}(n+1)u_{23}(n)+c_{43}(n+1)t_{33}(n)+c_{44}(n+1)t_{43}(n)\nonumber\\
\hphantom{\frac{\mathrm{d}}{\mathrm{d}\tau}t_{43}(n)=}{}
	-s_{42}(n)d_{23}(n)-t_{43}(n)c_{33}(n) . \label{eq-3.16}
\end{gather}
Here fourteen functions specifying the matrix elements $A_{jk}(n|z)$ are fixed by formulas
\begin{gather}\label{eq-3.17}
	a_{22}(n)=a_{22},
\\ \label{eq-3.18}
	c_{21}(n)=a_{22}t_{21}(n)/r_{22}(n),
\\ \label{eq-3.19}
	c_{12}(n+1)=t_{12}(n)a_{22}/r_{22}(n),
\\ \label{eq-3.20}
	b_{23}(n)=a_{22}s_{23}(n)/r_{22}(n),
\\ \label{eq-3.21}
	b_{32}(n+1)=s_{32}(n)a_{22}/r_{22}(n),
\\ \label{eq-3.22}
	b_{24}(n)=a_{22}s_{24}(n)/r_{22}(n),
\\ \label{eq-3.23}
	b_{42}(n+1)=s_{42}(n)a_{22}/r_{22}(n),
\\ \label{eq-3.24}
	e_{33}(n)=e_{33},
\\ \label{eq-3.25}
	c_{34}(n)=e_{33}t_{34}(n)/v_{33}(n),
\\ \label{eq-3.26}
	c_{43}(n+1)=t_{43}(n)e_{33}/v_{33}(n),
\\ \label{eq-3.27}
	d_{32}(n)=e_{33}u_{32}(n)/v_{33}(n),
\\ \label{eq-3.28}
	d_{23}(n+1)=u_{23}(n)e_{33}/v_{33}(n),
\\ \label{eq-3.29}
	d_{31}(n)=e_{33}u_{31}(n)/v_{33}(n),
\\ \label{eq-3.30}
	d_{13}(n+1)=u_{13}(n)e_{33}/v_{33}(n) ,
\end{gather}
whilst the rest four functions $c_{11}(n)$, $c_{22}(n)$, $c_{33}(n)$, $c_{44}(n)$ remain as yet being unfixed.

The above written sixteen semi-discrete equations (\ref{eq-3.1})--(\ref{eq-3.16}) and fourteen specification formulas (\ref{eq-3.17})--(\ref{eq-3.30}) are obtainable by the simple algebraic manipulations involving the zero-curvature equation (\ref{eq-2.1}) and the ans\"{a}tze (\ref{eq-2.2}) and (\ref{eq-2.3}) for the auxiliary matrices.

In general, the spatially independent parameters $a_{22}$ and $e_{33}$ can be arbitrary functions of time. Thus, the obtained set of semi-discrete nonlinear equations (\ref{eq-3.1})--(\ref{eq-3.16}) deciphered by the shorthand formulas (\ref{eq-3.17})--(\ref{eq-3.30}) can in general be the parametrically driven one. Such an evident but physically important property potentially admissible for a wide class of semi-discrete nonlinear integrable systems is usually overlooked or ignored by the scientific community.

\section{Constructive part of the on-site local conservation laws}
%\label{sec4}


The most constructive way to fix four arbitrary functions $c_{11}(n)$, $c_{22}(n)$, $c_{33}(n)$, $c_{44}(n)$ is based on the use of lowest local conservation laws
\begin{gather}\label{eq-4.1}
 \frac{\mathrm{d}}{\mathrm{d}\tau}\rho_{11}(n)=c_{11}(n+1)+c_{22}(n+1)+c_{33}(n+1) -c_{11}(n)-c_{22}(n)-c_{33}(n),
\\
 \frac{\mathrm{d}}{\mathrm{d}\tau}\rho_{22}(n)=c_{22}(n+1) + a_{22}\frac{s_{23}(n+1)s_{32}(n)}{r_{22}(n+1)r_{22}(n)} + a_{22}\frac{s_{24}(n+1)s_{42}(n)}{r_{22}(n+1)r_{22}(n)}\nonumber\\
 \hphantom{\frac{\mathrm{d}}{\mathrm{d}\tau}\rho_{22}(n)=}{}
 -c_{22}(n)-a_{22}\frac{s_{23}(n)s_{32}(n-1)}{r_{22}(n)r_{22}(n-1)}- a_{22}\frac{s_{24}(n)s_{42}(n-1)}{r_{22}(n)r_{22}(n-1)}, \label{eq-4.2}
\\
 \frac{\mathrm{d}}{\mathrm{d}\tau}\rho(n)=c_{11}(n+1)+c_{22}(n+1)+c_{33}(n+1)+c_{44}(n+1)\nonumber\\
 \hphantom{\frac{\mathrm{d}}{\mathrm{d}\tau}\rho(n)=}{}
 -c_{11}(n)-c_{22}(n)-c_{33}(n)-c_{44}(n),\label{eq-4.3}
\\
 \frac{\mathrm{d}}{\mathrm{d}\tau}\rho_{33}(n)=c_{33}(n+1) + e_{33}\frac{u_{32}(n+1)u_{23}(n)}{v_{33}(n+1)v_{33}(n)} + e_{33}\frac{u_{31}(n+1)u_{13}(n)}{v_{33}(n+1)v_{33}(n)}\nonumber\\
 \hphantom{\frac{\mathrm{d}}{\mathrm{d}\tau}\rho_{33}(n)=}{}
 -c_{33}(n)-e_{33}\frac{u_{32}(n)u_{23}(n-1)}{v_{33}(n)v_{33}(n-1)}- e_{33}\frac{u_{31}(n)u_{13}(n-1)}{v_{33}(n)v_{33}(n-1)}, \label{eq-4.4}
\\
 \frac{\mathrm{d}}{\mathrm{d}\tau}\rho_{44}(n)=c_{22}(n+1)+c_{33}(n+1)+c_{44}(n+1) -c_{22}(n)-c_{33}(n)-c_{44}(n) \label{eq-4.5}
\end{gather}
associated with the respective on-site local conserved densities
\begin{gather}
 \rho_{11}(n)=\ln[t_{12}(n)v_{33}(n)t_{21}(n)+u_{13}(n)t_{22}(n)u_{31}(n)\nonumber\\
 \hphantom{\rho_{11}(n)=\ln[}{}
	-t_{12}(n)u_{23}(n)u_{31}(n)-u_{13}(n)u_{32}(n)t_{21}(n)], \label{eq-4.6}
\\
	\rho_{22}(n)=\ln[r_{22}(n)], \label{eq-4.7}
\\
	\rho(n)=\ln[u_{13}(n)s_{42}(n)-t_{12}(n)t_{43}(n)]+\ln[u_{31}(n)s_{24}(n)-t_{21}(n)t_{34}(n)], \label{eq-4.8}
\\
	\rho_{33}(n)=\ln[v_{33}(n)], \label{eq-4.9}
\\
	\rho_{44}(n)=\ln[t_{43}(n)r_{22}(n)t_{34}(n)+s_{42}(n)t_{33}(n)s_{24}(n)\nonumber\\
\hphantom{\rho_{44}(n)=\ln[}{}
	-t_{43}(n)s_{32}(n)s_{24}(n)-s_{42}(n)s_{23}(n)t_{34}(n)] . \label{eq-4.10}
\end{gather}
These ten formulas (\ref{eq-4.1})--(\ref{eq-4.10}) are obtainable mainly within the modified direct recursive procedure \cite{vakhnenko-JNMP-18-401, vakhnenko-JNMP-18-415, vakhnenko-JNMP-20-606, vakhnenko-JMP-56-033505} originated as the generalization of Tsuchida--Ujino--Wadati direct recursive approach \cite{tsuchida-JMP-39-4785, tsuchida-JPA-32-2239}. Only two of them, namely formulas (\ref{eq-4.3}) and (\ref{eq-4.8}), arise as a simple paraphrase of so-called universal local conservation law
\begin{equation*}%\label{eq-4.11}
	\frac{\mathrm{d}}{\mathrm{d}\tau}\ln\left[\det L(n|z)\right]=\mathrm{Sp} A(n+1|z) -\mathrm{Sp} A(n|z)
\end{equation*}
following directly from the zero-curvature equation (\ref{eq-2.1}).

In order to fix four arbitrary functions $c_{11}(n)$, $c_{22}(n)$, $c_{33}(n)$, $c_{44}(n)$
we must impose four constraints onto the left-hand sides of five on-site local conservation laws (\ref{eq-4.1})--(\ref{eq-4.5}). Of course, such a procedure admits a number of diverse realizations giving rise to one or another particular sample of semi-discrete nonlinear integrable system. For this reason, the functions $c_{11}(n)$, $c_{22}(n)$, $c_{33}(n)$, $c_{44}(n)$ can be referred to as the sampling ones.

Meanwhile, our previous papers \cite{vakhnenko-JNMP-18-401, vakhnenko-JNMP-18-415} have substantially restricted the range of claimed diversity by ignoring the functions $c_{11}(n)$ and $c_{44}(n)$ in ansatz for the evolutionary matrix $A(n|z)$. In our present consideration, this ignorance is tantamount to the following two constrains:
\begin{gather*}%\label{eq-4.12}
\frac{\mathrm{d}}{\mathrm{d}\tau}\rho_{11}(n)=\frac{\mathrm{d}}{\mathrm{d}\tau}\rho(n) =\frac{\mathrm{d}}{\mathrm{d}\tau}\rho_{44}(n) .
\end{gather*}
The variability of another two admissible constrains has been thoroughly analyzed in the second~\cite{vakhnenko-JNMP-18-415} of just mentioned papers.

In what follows we try to grasp the key features of a particular semi-discrete nonlinear integrable system specified by four the most natural demands
\begin{gather}\label{eq-4.13}
	\frac{\mathrm{d}}{\mathrm{d}\tau}\rho_{22}(n)=0=\frac{\mathrm{d}}{\mathrm{d}\tau}\rho_{33}(n),
\\ \label{eq-4.14}
	\frac{\mathrm{d}}{\mathrm{d}\tau}\rho_{11}(n)=0=\frac{\mathrm{d}}{\mathrm{d}\tau}\rho_{44}(n) .
\end{gather}
Then the local conservation laws (\ref{eq-4.1}), (\ref{eq-4.2}) and (\ref{eq-4.4}), (\ref{eq-4.5}) yield
\begin{gather}
	c_{11}(n)=c_{11}+a_{22}\frac{s_{23}(n)s_{32}(n-1)}{r_{22}(n)r_{22}(n-1)}+ a_{22}\frac{s_{24}(n)s_{42}(n-1)}{r_{22}(n)r_{22}(n-1)}\nonumber\\
\hphantom{c_{11}(n)=}{}
+e_{33}\frac{u_{32}(n)u_{23}(n-1)}{v_{33}(n)v_{33}(n-1)}+ e_{33}\frac{u_{31}(n)u_{13}(n-1)}{v_{33}(n)v_{33}(n-1)},\label{eq-4.15}
\\
	c_{22}(n)=c_{22}-a_{22}\frac{s_{23}(n)s_{32}(n-1)}{r_{22}(n)r_{22}(n-1)}- a_{22}\frac{s_{24}(n)s_{42}(n-1)}{r_{22}(n)r_{22}(n-1)}, \label{eq-4.16}
\\
 c_{33}(n)=c_{33}-e_{33}\frac{u_{32}(n)u_{23}(n-1)}{v_{33}(n)v_{33}(n-1)}- e_{33}\frac{u_{31}(n)u_{13}(n-1)}{v_{33}(n)v_{33}(n-1)}, \label{eq-4.17}
\\
	c_{44}(n)=c_{44}+e_{33}\frac{u_{32}(n)u_{23}(n-1)}{v_{33}(n)v_{33}(n-1)}+ e_{33}\frac{u_{31}(n)u_{13}(n-1)}{v_{33}(n)v_{33}(n-1)}\nonumber\\
\hphantom{c_{44}(n)=}{}
	+a_{22}\frac{s_{23}(n)s_{32}(n-1)}{r_{22}(n)r_{22}(n-1)}+ a_{22}\frac{s_{24}(n)s_{42}(n-1)}{r_{22}(n)r_{22}(n-1)} . \label{eq-4.18}
\end{gather}
Here the time-dependent free parameters $c_{11}$, $c_{22}$, $c_{33}$, $c_{44}$ can be safely removed by the proper gauge transformations of field functions. Therefore, without the loss of generality we equalize each of four parameters $c_{jj}$ to zero.

As a result of adopted constraints (\ref{eq-4.13}) and (\ref{eq-4.14}), the number of independent field functions is reduced from the sixteen to the twelve ones. The choice of functions $t_{12}(n)$, $u_{13}(n)$, $t_{21}(n)$, $s_{23}(n)$, $u_{23}(n)$, $s_{24}(n)$, $u_{31}(n)$, $s_{32}(n)$, $u_{32}(n)$, $t_{34}(n)$, $s_{42}(n)$, $t_{43}(n)$ as being truly independent appears to be the most convenient. Thus, the functions $r_{22}(n)$, $t_{22}(n)$ and $v_{33}(n)$, $t_{33}(n)$ acquire the status of concomitant functions.

The arbitrary spatial dependencies of time independent concomitant functions~$r_{22}(n)$ and $v_{33}(n)$ could in principle imitate the action of external substrate, but here we discard this idea in order to preserve the uniformity of space. Thus, the concomitant functions $r_{22}(n)$ and $v_{33}(n)$ are reduced to the sheer constant parameters. The subsequent proper scaling procedure of field functions and the coupling parameters $a_{22}$ and $e_{33}$ gives rise to the equations of motion invariant to the primary equations of motion specified by the equalities
\begin{gather}\label{eq-4.19}
	r_{22}(n)=1=v_{33}(n) .
\end{gather}

As to the concomitant functions $t_{22}(n)$ and $t_{33}(n)$, they are determined by the simple algebraic equations
\begin{gather}
 t_{12}(n)v_{33}(n)t_{21}(n)+u_{13}(n)t_{22}(n)u_{31}(n)\nonumber\\
 \qquad{} -t_{12}(n)u_{23}(n)u_{31}(n)-u_{13}(n)u_{32}(n)t_{21}(n) = T_{22}(n),\label{eq-4.20}
\\
	t_{43}(n)r_{22}(n)t_{34}(n)+s_{42}(n)t_{33}(n)s_{24}(n)\nonumber\\
	\qquad{} -t_{43}(n)s_{32}(n)s_{24}(n)-s_{42}(n)s_{23}(n)t_{34}(n) = T_{33}(n)\label{eq-4.21}
\end{gather}
in accordance with the adopted differential constraints (\ref{eq-4.14}) accompanied by the explicit expressions (\ref{eq-4.6}) and (\ref{eq-4.10}) for the on-site local conserved densities $\rho_{11}(n)$ and $\rho_{44}(n)$.
Here $T_{22}(n)$ and $T_{22}(n)$ are time-independent arbitrary functions of spatial variable $n$.

In the present research, we restrict ourselves to the simplest possible variant
\begin{gather}\label{eq-4.22a}
	T_{22}(n) = 0 = T_{33}(n)
\end{gather}
supplemented by the earlier adopted normalizations (\ref{eq-4.19}) for $r_{22}(n)$ and $v_{33}(n)$. In addition, we introduce the transformation formulas
\begin{gather}\label{eq-4.22}
	t_{12}(n)=u_{13}(n)\bar{u}_{32}(n),
\\ \label{eq-4.24}
	t_{21}(n)=\bar{u}_{23}(n)u_{31}(n),
\\ \label{eq-4.25}
	t_{43}(n)=s_{42}(n)\bar{s}_{23}(n),
\\ \label{eq-4.26}
	t_{34}(n)=\bar{s}_{32}(n)s_{24}(n)
\end{gather}
serving to replace the field functions $t_{12}(n)$, $t_{21}(n)$, $t_{43}(n)$, $t_{34}(n)$ by the more suitable ones $\bar{u}_{32}(n)$, $\bar{u}_{23}(n)$, $\bar{s}_{23}(n)$, $\bar{s}_{32}(n)$. Then, the equations (\ref{eq-4.20}) and (\ref{eq-4.21}) for the concomitant functions $t_{22}(n)$ and $t_{33}(n)$ yield very simple results
\begin{gather*}%\label{eq-4.27}
	t_{22}(n)=u_{23}(n)\bar{u}_{32}(n)+\bar{u}_{23}(n)u_{32}(n)-\bar{u}_{23}(n)\bar{u}_{32}(n),
\\ %\label{eq-4.28}
	t_{33}(n)=s_{23}(n)\bar{s}_{32}(n) + \bar{s}_{23}(n)s_{32}(n) - \bar{s}_{23}(n)\bar{s}_{32}(n) .
\end{gather*}

We finalize this section by presenting four important local conserved densities \cite{vakhnenko-JNMP-18-401, vakhnenko-JNMP-18-415}
\begin{gather}\label{eq-4.29}
	\rho^{+}_{22}(n)=	\frac{t_{22}(n)}{r_{22}(n)} + \frac{s_{23}(n+1)s_{32}(n)}{r_{22}(n+1)r_{22}(n)} + \frac{s_{24}(n+1)s_{42}(n)}{r_{22}(n+1)r_{22}(n)},
\\ \label{eq-4.30}
	\rho^{-}_{22}(n)=	\frac{t_{22}(n)}{r_{22}(n)} + \frac{s_{23}(n)s_{32}(n-1)}{r_{22}(n)r_{22}(n-1)} + \frac{s_{24}(n)s_{42}(n-1)}{r_{22}(n)r_{22}(n-1)},
\\ \label{eq-4.31}
	\rho^{+}_{33}(n)=	\frac{t_{33}(n)}{v_{33}(n)} + \frac{u_{32}(n+1)u_{23}(n)}{v_{33}(n+1)v_{33}(n)} + \frac{u_{31}(n+1)u_{13}(n)}{v_{33}(n+1)v_{33}(n)},
\\ \label{eq-4.32}
	\rho^{-}_{33}(n)=	\frac{t_{33}(n)}{v_{33}(n)} + \frac{u_{32}(n)u_{23}(n-1)}{v_{33}(n)v_{33}(n-1)} + \frac{u_{31}(n)u_{13}(n-1)}{v_{33}(n)v_{33}(n-1)}
\end{gather}
that can be useful in constructing the appropriate Hamiltonian function for one or another particular realization of semi-discrete nonlinear integrable system associated with the adopted ans\"{a}tze (\ref{eq-2.2}) and (\ref{eq-2.3}) for the auxiliary spectral and evolutionary matrices.

The above listed local conserved densities (\ref{eq-4.29})--(\ref{eq-4.32}) are written in the most general form admitting any feasible fixation of sampling functions $c_{jj}(n)$.


\section[Appropriate reductions of field functions and coupling parameters]{Appropriate reductions of field functions\\ and coupling parameters}
\label{sec5}


Now all preliminary preparations have been completed and we are able to make appropriate reductions in the prototype set of semi-discrete equations (\ref{eq-3.1})--(\ref{eq-3.16}) by having taken into account the specification formulas (\ref{eq-3.17})--(\ref{eq-3.30}), (\ref{eq-4.15})--(\ref{eq-4.18}) for the functional elements of evolutionary matrix, the transformation formulas (\ref{eq-4.22})--(\ref{eq-4.26}) for more suitable field functions, as well as the expressions (\ref{eq-4.19})--(\ref{eq-4.22a}) for the concomitant quantities.

The proper consideration gives rise to the two types of reductions characterized by the twelve and six actual field functions, respectively.
In both of the announced reduction procedures, the parameter $\sigma$ is set to distinguish two admissible types of nonlinearities, namely, the attractive (focusing) $\sigma=+1$ and repulsive (defocusing) $\sigma=-1$ ones.

\subsection{Reduction to twelve field functions and two coupling parameters}\label{sec5.1}
The first type of reductions is specified by the following formulas:
\begin{alignat}{3}\label{eq-5.1}
	&s_{23}(n)=+q_{+}(n), \qquad&& u_{32}(n)=-\sigma r_{+}(n),&
\\ \label{eq-5.2}
	&u_{23}(n)=-q_{-}(n), \qquad&& s_{32}(n)=+\sigma r_{-}(n),&
\\ \label{eq-5.3}
	&\bar{s}_{23}(n)=+\bar{q}_{+}(n), \qquad&& \bar{u}_{32}(n)=-\sigma \bar{r}_{+}(n),&
\\ \label{eq-5.4}
	&\bar{u}_{23}(n)=-\bar{q}_{-}(n), \qquad&& \bar{s}_{32}(n)=+\sigma \bar{r}_{-}(n),&
\\ \label{eq-5.5}
	&s_{24}(n)=+f_{+}(n), \qquad&& u_{31}(n)=-\sigma g_{+}(n),&
\\ \label{eq-5.6}
	&u_{13}(n)=-f_{-}(n), \qquad&& s_{42}(n)=+\sigma g_{-}(n),&
\\ \label{eq-5.7}
	&t_{22}(n)=+\sigma \mu(n), \qquad&& t_{33}(n)=+\sigma \nu(n),&
\\ \label{eq-5.8}
	&a_{22}=-\mathrm{i}\alpha, \qquad&& e_{33}=+\mathrm{i}\beta .&
\end{alignat}

The twelve-component semi-discrete nonlinear integrable system written in terms of above introduced quantities (\ref{eq-5.1})--(\ref{eq-5.8}) is presented in Section~\ref{sec6}.

\subsection{Reduction to six field functions and one coupling parameter}\label{sec5.2}
The second type of reductions is specified by the following formulas:
\begin{alignat}{3}\label{eq-5.9}
&s_{23}(n)=+w_{+}(n), \qquad&& u_{32}(n)=-\sigma w_{+}(n),&
\\ \label{eq-5.10}
&u_{23}(n)=-w_{-}(n), \qquad&& s_{32}(n)=+\sigma w_{-}(n),&
\\ \label{eq-5.11}
&\bar{s}_{23}(n)=+\bar{w}_{+}(n), \qquad&& \bar{u}_{32}(n)=-\sigma \bar{w}_{+}(n),&
\\ \label{eq-5.12}
&\bar{u}_{23}(n)=-\bar{w}_{-}(n), \qquad&& \bar{s}_{32}(n)=+\sigma \bar{w}_{-}(n),&
\\ \label{eq-5.13}
&s_{24}(n)=+h_{+}(n), \qquad&& u_{31}(n)=-\sigma h_{+}(n),&
\\ \label{eq-5.14}
&u_{13}(n)=-h_{-}(n), \qquad&& s_{42}(n)=+\sigma h_{-}(n),&
\\ \label{eq-5.15}
&t_{22}(n)=+\sigma \pi(n), \qquad&& t_{33}(n)=+\sigma \pi(n),&
\\ \label{eq-5.16}
&a_{22}=\varkappa, \qquad&& e_{33}(n)=\varkappa .&
\end{alignat}

The six-component semi-discrete nonlinear integrable system written in terms of above introduced quantities (\ref{eq-5.9})--(\ref{eq-5.16}) is presented in Section~\ref{sec7}.


\section[Twelve-component semi-discrete nonlinear integrable system and its admissible symmetries]{Twelve-component semi-discrete nonlinear integrable system\\ and its admissible symmetries}
\label{sec6}


The reduction formulas (\ref{eq-5.1})--(\ref{eq-5.8}) listed in Section~\ref{sec5.1} as applied to the prototype set of semi-discrete equations (\ref{eq-3.1})--(\ref{eq-3.16})
accompanied by the formulas referred in the first paragraph of Section~\ref{sec5} give rise to the following twelve-component semi-discrete nonlinear integrable system:
\begin{gather}
	 +\mathrm{i}\frac{\mathrm{d}}{\mathrm{d}\tau}q_{+}(n) = - \alpha q_{-}(n) - \beta q_{-}(n-1)[1+\sigma r_{+}(n)q_{+}(n)] - \alpha\sigma\mu(n)q_{+}(n)\nonumber\\
 \hspace{0.25mm}\hphantom{+\mathrm{i}\frac{\mathrm{d}}{\mathrm{d}\tau}q_{+}(n) =}{}
	 +\alpha\sigma q_{+}(n+1)[\nu(n)-r_{-}(n)q_{+}(n)] - \alpha\sigma f_{+}(n+1)g_{-}(n)[q_{+}(n)-\bar{q}_{+}(n)]\nonumber\\
 \hspace{0.25mm}\hphantom{+\mathrm{i}\frac{\mathrm{d}}{\mathrm{d}\tau}q_{+}(n) =}{}
	 -\beta\sigma f_{-}(n-1)g_{+}(n)q_{+}(n) +\beta\sigma f_{+}(n)g_{-}(n-1)\bar{q}_{+}(n-1), \label{eq-6.1}
\\
	 -\mathrm{i}\frac{\mathrm{d}}{\mathrm{d}\tau}r_{+}(n) = -\beta r_{-}(n) - \alpha r_{-}(n-1)[1+\sigma q_{+}(n)r_{+}(n)] - \beta\sigma\nu(n)r_{+}(n)\nonumber\\
 \hspace{0.25mm}\hphantom{-\mathrm{i}\frac{\mathrm{d}}{\mathrm{d}\tau}r_{+}(n) =}{}
	 +\beta\sigma r_{+}(n+1)[\mu(n)-q_{-}(n)r_{+}(n)] - \beta\sigma g_{+}(n+1)f_{-}(n)[r_{+}(n)-\bar{r}_{+}(n)]\nonumber\\
 \hspace{0.25mm}\hphantom{-\mathrm{i}\frac{\mathrm{d}}{\mathrm{d}\tau}r_{+}(n) =}{}
	 -\alpha\sigma g_{-}(n-1)f_{+}(n)r_{+}(n) +\alpha\sigma g_{+}(n)f_{-}(n-1)\bar{r}_{+}(n-1), \label{eq-6.2}
\\
	 +\mathrm{i}\frac{\mathrm{d}}{\mathrm{d}\tau}q_{-}(n) = -\beta q_{+}(n) - \alpha q_{+}(n+1)[1+\sigma r_{-}(n)q_{-}(n)] - \beta\sigma\nu(n)q_{-}(n)\nonumber\\
 \hspace{0.25mm}\hphantom{+\mathrm{i}\frac{\mathrm{d}}{\mathrm{d}\tau}q_{-}(n) =}{}
	 +\beta\sigma q_{-}(n-1)[\mu(n)-r_{+}(n)q_{-}(n)] - \beta\sigma f_{-}(n-1)g_{+}(n)[q_{-}(n)-\bar{q}_{-}(n)]\nonumber\\
 \hspace{0.25mm}\hphantom{+\mathrm{i}\frac{\mathrm{d}}{\mathrm{d}\tau}q_{-}(n) =}{}
	 -\alpha\sigma f_{+}(n+1)g_{-}(n)q_{-}(n) +\alpha\sigma f_{-}(n)g_{+}(n+1)\bar{q}_{-}(n+1), \label{eq-6.3}
\\
	 -\mathrm{i}\frac{\mathrm{d}}{\mathrm{d}\tau}r_{-}(n) = -\alpha r_{+}(n) - \beta r_{+}(n+1)[1+\sigma q_{-}(n)r_{-}(n)] - \alpha\sigma\mu(n)r_{-}(n)\nonumber\\
 \hspace{0.25mm}\hphantom{-\mathrm{i}\frac{\mathrm{d}}{\mathrm{d}\tau}r_{-}(n) =}{}
	 +\alpha\sigma r_{-}(n-1)[\nu(n)-q_{+}(n)r_{-}(n)] - \alpha\sigma g_{-}(n-1)f_{+}(n)[r_{-}(n)-\bar{r}_{-}(n)]\nonumber\\
 \hspace{0.25mm}\hphantom{-\mathrm{i}\frac{\mathrm{d}}{\mathrm{d}\tau}r_{-}(n) =}{}
	 -\beta\sigma g_{+}(n+1)f_{-}(n)r_{-}(n) +\beta\sigma g_{-}(n)f_{+}(n+1)\bar{r}_{-}(n+1), \label{eq-6.4}
\\
	 +\mathrm{i}\frac{\mathrm{d}}{\mathrm{d}\tau}\bar{q}_{+}(n) = -\alpha q_{-}(n) - \beta q_{-}(n-1)[1+\sigma r_{+}(n)\bar{q}_{+}(n)] - \alpha\sigma\mu(n)\bar{q}_{+}(n)\nonumber\\
 \hspace{0.25mm}\hphantom{+\mathrm{i}\frac{\mathrm{d}}{\mathrm{d}\tau}\bar{q}_{+}(n) =}{}
	 -\beta\sigma \bar{q}_{+}(n)[\nu(n)-r_{-}(n)\bar{q}_{+}(n)] - \alpha\sigma \bar{q}_{+}(n)r_{-}(n-1)[q_{+}(n)-\bar{q}_{+}(n)]\nonumber\\
 \hspace{0.25mm}\hphantom{+\mathrm{i}\frac{\mathrm{d}}{\mathrm{d}\tau}\bar{q}_{+}(n) =}{}
	 -\alpha\sigma f_{+}(n)g_{-}(n-1)\bar{q}_{+}(n) -\beta\sigma f_{-}(n-1)g_{+}(n)\bar{q}_{+}(n), \label{eq-6.5}
\\
	 -\mathrm{i}\frac{\mathrm{d}}{\mathrm{d}\tau}\bar{r}_{+}(n) = -\beta r_{-}(n) - \alpha r_{-}(n-1)[1+\sigma q_{+}(n)\bar{r}_{+}(n)] - \beta\sigma\nu(n)\bar{r}_{+}(n)\nonumber\\
 \hspace{0.25mm}\hphantom{-\mathrm{i}\frac{\mathrm{d}}{\mathrm{d}\tau}\bar{r}_{+}(n) =}{}
	 -\alpha\sigma \bar{r}_{+}(n)[\mu(n)-q_{-}(n)\bar{r}_{+}(n)] - \beta\sigma \bar{r}_{+}(n)q_{-}(n-1)[r_{+}(n)-\bar{r}_{+}(n)]\nonumber\\
 \hspace{0.25mm}\hphantom{-\mathrm{i}\frac{\mathrm{d}}{\mathrm{d}\tau}\bar{r}_{+}(n) =}{}
	 -\beta\sigma g_{+}(n)f_{-}(n-1)\bar{r}_{+}(n) -\alpha\sigma g_{-}(n-1)f_{+}(n)\bar{r}_{+}(n), \label{eq-6.6}
\\
	 +\mathrm{i}\frac{\mathrm{d}}{\mathrm{d}\tau}\bar{q}_{-}(n) = -\beta q_{+}(n) - \alpha q_{+}(n+1)[1+\sigma r_{-}(n)\bar{q}_{-}(n)] - \beta\sigma\nu(n)\bar{q}_{-}(n)\nonumber\\
 \hspace{0.25mm}\hphantom{+\mathrm{i}\frac{\mathrm{d}}{\mathrm{d}\tau}\bar{q}_{-}(n) =}{}
	 -\alpha\sigma \bar{q}_{-}(n)[\mu(n)-r_{+}(n)\bar{q}_{-}(n)] - \beta\sigma \bar{q}_{-}(n)r_{+}(n+1)[q_{-}(n)-\bar{q}_{-}(n)]\nonumber\\
 \hspace{0.25mm}\hphantom{+\mathrm{i}\frac{\mathrm{d}}{\mathrm{d}\tau}\bar{q}_{-}(n) =}{}
	 -\beta\sigma f_{-}(n)g_{+}(n+1)\bar{q}_{-}(n) -\alpha\sigma f_{+}(n+1)g_{-}(n)\bar{q}_{-}(n), \label{eq-6.7}
\\
	 -\mathrm{i}\frac{\mathrm{d}}{\mathrm{d}\tau}\bar{r}_{-}(n) = -\alpha r_{+}(n) - \beta r_{+}(n+1)[1+\sigma q_{-}(n)\bar{r}_{-}(n)] - \alpha\sigma\mu(n)\bar{r}_{-}(n)\nonumber\\
 \hspace{0.25mm}\hphantom{-\mathrm{i}\frac{\mathrm{d}}{\mathrm{d}\tau}\bar{r}_{-}(n) =}{}
	 -\beta\sigma \bar{r}_{-}(n)[\nu(n)-q_{+}(n)\bar{r}_{-}(n)] - \alpha\sigma \bar{r}_{-}(n)q_{+}(n+1)[r_{-}(n)-\bar{r}_{-}(n)]\nonumber\\
 \hspace{0.25mm}\hphantom{-\mathrm{i}\frac{\mathrm{d}}{\mathrm{d}\tau}\bar{r}_{-}(n) =}{}
	 -\alpha\sigma g_{-}(n)f_{+}(n+1)\bar{r}_{-}(n) -\beta\sigma g_{+}(n+1)f_{-}(n)\bar{r}_{-}(n), \label{eq-6.8}
\\
	 +\mathrm{i}\sigma\frac{\mathrm{d}}{\mathrm{d}\tau}\ln[f_{+}(n)] = -\alpha \mu(n) + \beta q_{+}(n)\bar{r}_{-}(n)\nonumber\\
 \hspace{0.25mm}\hphantom{+\mathrm{i}\sigma\frac{\mathrm{d}}{\mathrm{d}\tau}\ln[f_{+}(n)] =}{}
	 + \beta q_{-}(n-1)r_{+}(n) - \alpha q_{+}(n)r_{-}(n-1) - \alpha q_{+}(n+1)[r_{-}(n)-\bar{r}_{-}(n)]\nonumber\\
 \hspace{0.25mm}\hphantom{+\mathrm{i}\sigma\frac{\mathrm{d}}{\mathrm{d}\tau}\ln[f_{+}(n)] =}{}
	 - \alpha f_{+}(n+1)g_{-}(n) - \alpha f_{+}(n)g_{-}(n-1) + \beta f_{-}(n-1)g_{+}(n), \label{eq-6.9}
\\
	 -\mathrm{i}\sigma\frac{\mathrm{d}}{\mathrm{d}\tau}\ln[g_{+}(n)] = -\beta \nu(n) + \alpha r_{+}(n)\bar{q}_{-}(n)\nonumber\\
 \hspace{0.25mm}\hphantom{-\mathrm{i}\sigma\frac{\mathrm{d}}{\mathrm{d}\tau}\ln[g_{+}(n)] =}{}
	 + \alpha r_{-}(n-1)q_{+}(n) - \beta r_{+}(n)q_{-}(n-1) - \beta r_{+}(n+1)[q_{-}(n)-\bar{q}_{-}(n)]\nonumber\\
 \hspace{0.25mm}\hphantom{-\mathrm{i}\sigma\frac{\mathrm{d}}{\mathrm{d}\tau}\ln[g_{+}(n)] =}{}
	 - \beta g_{+}(n+1)f_{-}(n) - \beta g_{+}(n)f_{-}(n-1) + \alpha g_{-}(n-1)f_{+}(n), \label{eq-6.10}
\\
	 +\mathrm{i}\sigma\frac{\mathrm{d}}{\mathrm{d}\tau}\ln[f_{-}(n)] = -\beta \nu(n) + \alpha q_{-}(n)\bar{r}_{+}(n)\nonumber\\
 \hspace{0.25mm}\hphantom{+\mathrm{i}\sigma\frac{\mathrm{d}}{\mathrm{d}\tau}\ln[f_{-}(n)] =}{}
	 + \alpha q_{+}(n+1)r_{-}(n) - \beta q_{-}(n)r_{+}(n+1) - \beta q_{-}(n-1)[r_{+}(n)-\bar{r}_{+}(n)]\nonumber\\
 \hspace{0.25mm}\hphantom{+\mathrm{i}\sigma\frac{\mathrm{d}}{\mathrm{d}\tau}\ln[f_{-}(n)] =}{}
	 - \beta f_{-}(n-1)g_{+}(n) - \beta f_{-}(n)g_{+}(n+1) + \alpha f_{+}(n+1)g_{-}(n), \label{eq-6.11}
\\
	 -\mathrm{i}\sigma\frac{\mathrm{d}}{\mathrm{d}\tau}\ln[g_{-}(n)] = -\alpha \mu(n) + \beta r_{-}(n)\bar{q}_{+}(n)\nonumber\\
 \hspace{0.25mm}\hphantom{-\mathrm{i}\sigma\frac{\mathrm{d}}{\mathrm{d}\tau}\ln[g_{-}(n)] =}{}
	 + \beta r_{+}(n+1)q_{-}(n) - \alpha r_{-}(n)q_{+}(n+1) - \alpha r_{-}(n-1)[q_{+}(n)-\bar{q}_{+}(n)]\nonumber\\
 \hspace{0.25mm}\hphantom{-\mathrm{i}\sigma\frac{\mathrm{d}}{\mathrm{d}\tau}\ln[g_{-}(n)] =}{}
	 - \alpha g_{-}(n-1)f_{+}(n) - \alpha g_{-}(n)f_{+}(n+1) + \beta g_{+}(n+1)f_{-}(n) . \label{eq-6.12}
\end{gather}
Here the concomitant field functions $\mu(n)$ and $\nu(n)$ are given by formulas
\begin{gather}\label{eq-6.13}
	\mu(n) = q_{-}(n)\bar{r}_{+}(n)+\bar{q}_{-}(n)r_{+}(n)-\bar{q}_{-}(n)\bar{r}_{+}(n),
\\ \label{eq-6.14}
	\nu(n) = r_{-}(n)\bar{q}_{+}(n)+\bar{r}_{-}(n)q_{+}(n)-\bar{r}_{-}(n)\bar{q}_{+}(n) .
\end{gather}
We would like to remind that the spatially independent parameters $\alpha$ and $\beta$ can be arbitrary functions of time, while the parameter $\sigma$ defined as $\sigma^{2}=1$ serves to distinguish two types of nonlinearities.

It is worth noticing that the twelve-component system (\ref{eq-6.1})--(\ref{eq-6.14}) can be treated as settled on two mutually coupled chains marked by indices $+$ and $-$. In this sense, the lattice as a whole is proved to be a quasi-one-dimensional one. The system's spatial quasi-one-dimensionality and multicomponentness could be very prospective attributes in modeling transport properties of long macromolecules both natural and synthetic origins \cite{chen-AM-32-2001893, perepelytsya-EBJ-50-759, satchanska-PM-16-1159}.

The obtained dynamical system (\ref{eq-6.1})--(\ref{eq-6.14}) admits at least two types of symmetries.

\subsection{Symmetry under the complex conjugation}
Thus, the system's symmetry under the complex conjugation is based on the complex conjugate symmetries of field functions
\begin{gather}\label{eq-6.15}
	 r^{*}_{+}(n)=q_{+}(n),
\\ \label{eq-6.16}
	 r^{*}_{-}(n)=q_{-}(n),
\\ \label{eq-6.17}
	 \bar{r}^{*}_{+}(n)=\bar{q}_{+}(n),
\\ \label{eq-6.18}
	 \bar{r}^{*}_{-}(n)=\bar{q}_{-}(n),
\\ \label{eq-6.19}
	 g^{*}_{+}(n)=f_{+}(n),
\\ \label{eq-6.20}
	 g^{*}_{-}(n)=f_{-}(n),
\\ \label{eq-6.21}
	\nu^{*}(n)=\mu(n)
\end{gather}
that are valid provided the parameters $\alpha$, $-\mathrm{i}$, $\beta$
meet the conditions of complex conjugation{\samepage
\begin{gather}\label{eq-6.22}
	 \beta^{*}=\alpha,
\\ \label{eq-6.23}
	(-\mathrm{i})^{*}=+\mathrm{i} .
\end{gather}}%
Due to their complexvalueness, the coupling parameters $\alpha$ and $\beta$ are able to model the impact of external uniform magnetic field in terms of Peierls phase factors \cite{feynman-FLP3-1963, peierls-ZPh-80-763, vakhnenko-PRE-77-026604}.

\subsection{Symmetry under the space and time reversal}
The system's symmetry under the space and time reversal is more sophisticated and it is based on the dynamical properties of following transformed field functions
\begin{gather}\label{eq-6.24}
	\mathrm{q}_{+}(n)\equiv\mathrm{q}_{+}(n|\tau)=r_{-}(-n|-\tau),
\\ \label{eq-6.25}
	\mathrm{r}_{+}(n)\equiv\mathrm{r}_{+}(n|\tau)=q_{-}(-n|-\tau),
\\ \label{eq-6.26}
	\mathrm{q}_{-}(n)\equiv\mathrm{q}_{-}(n|\tau)=r_{+}(-n|-\tau),
\\ \label{eq-6.27}
	\mathrm{r}_{-}(n)\equiv\mathrm{r}_{-}(n|\tau)=q_{+}(-n|-\tau),
\\ \label{eq-6.28}
	\bar{\mathrm{q}}_{+}(n)\equiv\bar{\mathrm{q}}_{+}(n|\tau)=\bar{r}_{-}(-n|-\tau),
\\ \label{eq-6.29}
	\bar{\mathrm{r}}_{+}(n)\equiv\bar{\mathrm{r}}_{+}(n|\tau)=\bar{q}_{-}(-n|-\tau),
\\ \label{eq-6.30}
	\bar{\mathrm{q}}_{-}(n)\equiv\bar{\mathrm{q}}_{-}(n|\tau)=\bar{r}_{+}(-n|-\tau),
\\ \label{eq-6.31}
	\bar{\mathrm{r}}_{-}(n)\equiv\bar{\mathrm{r}}_{-}(n|\tau)=\bar{q}_{+}(-n|-\tau),
\\ \label{eq-6.32}
	\bar{\mathrm{f}}_{+}(n)\equiv\bar{\mathrm{f}}_{+}(n|\tau)=g_{-}(-n|-\tau),
\\ \label{eq-6.33}
	\bar{\mathrm{g}}_{+}(n)\equiv\bar{\mathrm{g}}_{+}(n|\tau)=f_{-}(-n|-\tau),
\\ \label{eq-6.34}
	\bar{\mathrm{f}}_{-}(n)\equiv\bar{\mathrm{f}}_{-}(n|\tau)=g_{+}(-n|-\tau),
\\ \label{eq-6.35}
	\bar{\mathrm{g}}_{-}(n)\equiv\bar{\mathrm{g}}_{-}(n|\tau)=f_{+}(-n|-\tau),
\\ %\label{eq-6.36}
\upmu (n)\equiv\upmu(n|\tau)=\mu(-n|-\tau), \nonumber
\\ %\label{eq-6.37}
	\upnu(n)\equiv\upnu(n|\tau)=\nu(-n|-\tau) . \nonumber
\end{gather}
The simple comparison shows that the transformed field functions (\ref{eq-6.24})--(\ref{eq-6.35}) are governed by the set of twelve semi-discrete nonlinear equations invariant to the twelve-component semi-discrete nonlinear integrable equations of our interest (\ref{eq-6.1})--(\ref{eq-6.12}).

Here any additional requirements on parameters $\alpha$, $\mathrm{i}$, $\beta$ are seen to be unnecessary. In this sense the space-time reversal symmetry of inspected semi-discrete nonlinear integrable system (\ref{eq-6.1})--(\ref{eq-6.12}) turns out to be more general than the usual parity-time ($\cal{PT}$) symmetry \cite{bender-RPP-70-947, konotop-RMP-88-035002}.


\section[Six-component semi-discrete nonlinear integrable system and its symmetry under the space and time reversal]{Six-component semi-discrete nonlinear integrable system \\
and its symmetry under the space and time reversal}
\label{sec7}


The reduction formulas (\ref{eq-5.9})--(\ref{eq-5.16}) listed in Section~\ref{sec5.2} as applied to the prototype set of semi-discrete equations (\ref{eq-3.1})--(\ref{eq-3.16})
accompanied by the formulas referred in the first paragraph of Section~\ref{sec5} give rise to the following six-component semi-discrete nonlinear integrable system:
\begin{gather}
	 \frac{\mathrm{d}}{\mathrm{d}\tau}w_{+}(n) = -\varkappa w_{-}(n) + \varkappa w_{-}(n-1)[1+\sigma w_{+}(n)w_{+}(n)] - \varkappa\sigma\pi(n)w_{+}(n)\nonumber\\
\hspace{0.25mm}\hphantom{\frac{\mathrm{d}}{\mathrm{d}\tau}w_{+}(n) =}{}
	 +\varkappa\sigma w_{+}(n+1)[\pi(n)-w_{-}(n)w_{+}(n)] - \varkappa\sigma h_{+}(n+1)h_{-}(n)[w_{+}(n)-\bar{w}_{+}(n)]\nonumber\\
\hspace{0.25mm}\hphantom{\frac{\mathrm{d}}{\mathrm{d}\tau}w_{+}(n) =}{}
	 +\varkappa\sigma h_{+}(n)h_{-}(n-1)[w_{+}(n) - \bar{w}_{+}(n-1)],\label{eq-7.1}
\\
	 \frac{\mathrm{d}}{\mathrm{d}\tau}w_{-}(n) = + \varkappa w_{+}(n) - \varkappa w_{+}(n+1)[1+\sigma w_{-}(n)w_{-}(n)] + \varkappa\sigma\pi(n)w_{-}(n)\nonumber\\
\hspace{0.25mm}\hphantom{\frac{\mathrm{d}}{\mathrm{d}\tau}w_{-}(n) =}{}
	 -\varkappa\sigma w_{-}(n-1)[\pi(n)-w_{+}(n)w_{-}(n)] + \varkappa\sigma h_{-}(n-1)h_{+}(n)[w_{-}(n)-\bar{w}_{-}(n)]\nonumber\\
\hspace{0.25mm}\hphantom{\frac{\mathrm{d}}{\mathrm{d}\tau}w_{-}(n) =}{}
	 -\varkappa\sigma h_{-}(n)h_{+}(n+1)[w_{-}(n) -\bar{w}_{-}(n+1)], \label{eq-7.2}
\\
	 \frac{\mathrm{d}}{\mathrm{d}\tau}\bar{w}_{+}(n) = -\varkappa w_{-}(n)[1+\sigma \bar{w}_{+}(n)\bar{w}_{+}(n)] + \varkappa w_{-}(n-1)[1+\sigma \bar{w}_{+}(n)\bar{w}_{+}(n)], \label{eq-7.3}
\\
	 \frac{\mathrm{d}}{\mathrm{d}\tau}\bar{w}_{-}(n) = +\varkappa w_{+}(n)[1+\sigma \bar{w}_{-}(n)\bar{w}_{-}(n)] - \varkappa w_{+}(n+1)[1+\sigma \bar{w}_{-}(n)\bar{w}_{-}(n)], \label{eq-7.4}
\\
	 \sigma\frac{\mathrm{d}}{\mathrm{d}\tau}\ln[h_{+}(n)] = -\varkappa \pi(n) - \varkappa w_{+}(n)\bar{w}_{-}(n) - 2\varkappa w_{-}(n-1)w_{+}(n) \label{eq-7.5}\\
\hspace{0.25mm}\hphantom{\sigma\frac{\mathrm{d}}{\mathrm{d}\tau}\ln[h_{+}(n)] =}{}
	 - \varkappa w_{+}(n\hspace{-0.4pt}+\hspace{-0.4pt}1)[w_{-}(n)\hspace{-0.4pt}-\hspace{-0.4pt}\bar{w}_{-}(n)]
	\hspace{-0.4pt}-\hspace{-0.4pt} \varkappa h_{+}(n+1)h_{-}(n)\hspace{-0.4pt} -\hspace{-0.4pt} 2\varkappa h_{-}(n\hspace{-0.4pt}-\hspace{-0.4pt}1)h_{+}(n),\nonumber
\\
	 \sigma\frac{\mathrm{d}}{\mathrm{d}\tau}\ln[h_{-}(n)] = +\varkappa \pi(n) + \varkappa w_{-}(n)\bar{w}_{+}(n) + 2\varkappa w_{+}(n+1)w_{-}(n) \label{eq-7.6}\\
\hspace{0.25mm}\hphantom{\sigma\frac{\mathrm{d}}{\mathrm{d}\tau}\ln[h_{-}(n)] =}{}
	 + \varkappa w_{-}(n\hspace{-0.4pt}-\hspace{-0.4pt}1)[w_{+}(n)\hspace{-0.4pt}-\hspace{-0.4pt}\bar{w}_{+}(n)]
	 + \varkappa h_{-}(n\hspace{-0.4pt}-\hspace{-0.4pt}1)h_{+}(n)\hspace{-0.4pt} +\hspace{-0.4pt} 2\varkappa h_{+}(n\hspace{-0.4pt}+\hspace{-0.4pt}1)h_{-}(n) .\nonumber
\end{gather}
Here the concomitant field function $\pi(n)$ is given by formula
\begin{gather*}%\label{eq-7.7}
	\pi(n)=w_{-}(n)\bar{w}_{+}(n)+\bar{w}_{-}(n)w_{+}(n)-\bar{w}_{-}(n)\bar{w}_{+}(n) .
\end{gather*}
The system's symmetry under the space and time reversal is based on the dynamical properties of following transformed field functions:
\begin{gather}\label{eq-7.8}
	\mathrm{w}_{+}(n)\equiv\mathrm{w}_{+}(n|\tau)=w_{-}(-n|-\tau),
\\ \label{eq-7.25}
	\mathrm{w}_{-}(n)\equiv\mathrm{w}_{-}(n|\tau)=w_{+}(-n|-\tau),
\\ \label{eq-7.26}
	\bar{\mathrm{w}}_{+}(n)\equiv\bar{\mathrm{w}}_{+}(n|\tau)=\bar{w}_{-}(-n|-\tau),
\\ \label{eq-7.27}
	\bar{\mathrm{w}}_{-}(n)\equiv\bar{\mathrm{w}}_{-}(n|\tau)=\bar{w}_{+}(-n|-\tau),
\\ \label{eq-7.28}
	\bar{\mathrm{h}}_{+}(n)\equiv\bar{\mathrm{h}}_{+}(n|\tau)=h_{-}(-n|-\tau),
\\ \label{eq-7.29}
	\bar{\mathrm{h}}_{-}(n)\equiv\bar{\mathrm{h}}_{-}(n|\tau)=h_{+}(-n|-\tau),
\\ %\label{eq-7.30}
	\uppi (n)\equiv\uppi(n|\tau)=\pi(-n|-\tau) . \nonumber
\end{gather}
The simple comparison shows that the transformed field functions (\ref{eq-7.8})--(\ref{eq-7.29}) are governed by the set of six semi-discrete nonlinear equations invariant to the six-component semi-discrete nonlinear integrable equations of our interest (\ref{eq-7.1})--(\ref{eq-7.6}).

\section{Discussion}
%\label{sec8}

The semi-discrete nonlinear integrable systems (\ref{eq-6.1})--(\ref{eq-6.12}) and (\ref{eq-7.1})--(\ref{eq-7.6}) presented in Sections~\ref{sec6} and~\ref{sec7} are proved to be essentially multicomponent ones inasmuch either of them cannot be split into several physically uncoupled subsystems.

Thus, the twelve-component semi-discrete integrable system (\ref{eq-6.1})--(\ref{eq-6.12}) is composed of six subsystems coupled by linear and nonlinear types of interactions. These subsystems are formalized by the six pairs of field functions.
The set of plausible pairs are as follows
${q_{+}(n)\leftrightarrow r_{+}(n)}$, $\bar{q}_{+}(n)\leftrightarrow \bar{r}_{+}(n)$, $f_{+}(n)\leftrightarrow g_{+}(n)$, $q_{-}(n)\leftrightarrow r_{-}(n)$, $\bar{q}_{-}(n)\leftrightarrow \bar{r}_{-}(n)$, $f_{-}(n)\leftrightarrow g_{-}(n)$. Here the symbol $\leftrightarrow$ is inserted to point out on a suppositional canonical relationship between the field functions of a particular pair. The problems of adequate Hamiltonian treatment and reliable Poisson structure characterizing the twelve-component semi-discrete nonlinear integrable system (\ref{eq-6.1})--(\ref{eq-6.12})
appear to be very complicated and presently they are opened for the future investigations. The general principles of Hamiltonian and Poisson approaches as applied to the multicomponent semi-discrete nonlinear integrable systems have been approbated in our previous papers \cite{vakhnenko-JMP-56-033505, vakhnenko-JMP-57-113504,vakhnenko-JNMP-24-250, vakhnenko-JMP-59-053504, vakhnenko-LMP-108-1807}. The difficult aspects in establishing the Hamiltonian and Poisson structures for the systems of present interest (\ref{eq-6.1})--(\ref{eq-6.12}) and (\ref{eq-7.1})--(\ref{eq-7.6}) are summarized in~Appendix~\ref{app}.

Though the Hamiltonian treatment of twelve-component semi-discrete integrable system (\ref{eq-6.1})--(\ref{eq-6.12}) is waiting for its rigorous formulation the analysis of universal local conservation law (\ref{eq-4.3}) provides us with certain fruitful information about plausible physical sense of involved subsystems at least in the case of system's complex conjugate symmetry.

To be precise, we should inspect the expressions for the local conserved density $\rho(n)$ (\ref{eq-4.8}) and the local current $J(n)$ having been adapted to the needs of reduced semi-discrete nonlinear integrable system (\ref{eq-6.1})--(\ref{eq-6.12}) under the premise of complex conjugation symmetry (\ref{eq-6.15})--(\ref{eq-6.23}). The announced adapted quantities as well as the respective local conservation law read as follows:
\begin{gather}%\label{eq-8.1}
	\rho(n) = \ln[1+\sigma\bar{q}_{+}(n)\bar{r}_{+}(n)] + \ln[1+\sigma\bar{q}_{-}(n)\bar{r}_{-}(n)] + \ln[f_{+}(n)g_{+}(n)f_{-}(n)g_{-}(n)],\nonumber
\\
	J(n)=\mathrm{i}\alpha\sigma q_{+}(n)r_{-}(n-1)+\mathrm{i}\alpha\sigma f_{+}(n)g_{-}(n-1)\nonumber\\
\hphantom{J(n)=}{}
 -\mathrm{i}\beta\sigma r_{+}(n)q_{-}(n-1) - \mathrm{i}\beta\sigma g_{+}(n)f_{-}(n-1),\label{eq-8.2}
\\ %\label{eq-8.3}
	\frac{\mathrm{d}}{\mathrm{d}\tau}\rho(n) = J(n) - J(n+1) .\nonumber
\end{gather}
Here the partial local densities
\begin{gather}\label{eq-8.4}
	\bar{\rho}_{+}(n) = \ln[1+\sigma\bar{q}_{+}(n)\bar{r}_{+}(n)],
\\
\label{eq-8.5}
	\bar{\rho}_{-}(n) = \ln[1+\sigma\bar{q}_{-}(n)\bar{r}_{-}(n)]
\end{gather}
are essentially separate characteristics related to subsystems
 $\bar{q}_{+}(n)\leftrightarrow \bar{r}_{+}(n)$ and $\bar{q}_{-}(n)\leftrightarrow \bar{r}_{-}(n)$, respectively.
On the other hand, the net local density
\begin{gather}\label{eq-8.6}
	\rho_{\pm}(n) = \ln[f_{+}(n)g_{+}(n)f_{-}(n)g_{-}(n)]
\end{gather}
is related to two subsystems $f_{+}(n)\leftrightarrow g_{+}(n)$ and $f_{-}(n) \leftrightarrow g_{-}(n)$ combined.

In the case of attractive nonlinearity $\sigma=+1$, the first two local densities (\ref{eq-8.4}) and (\ref{eq-8.5}) are seen to be the real-valued nonnegative quantities treatable as the local densities of positive charges associated with the respective fields $\bar{q}_{+}(n)\leftrightarrow \bar{r}_{+}(n)$ and $\bar{q}_{-}(n)\leftrightarrow \bar{r}_{-}(n)$. In contrast, the third real-valued quantity (\ref{eq-8.6}) can acquire either positive or negative magnitude. This quantity can be treated as the net local density of charge related to two subsystems described by two pairs of fields $f_{+}(n)\leftrightarrow g_{+}(n)$ and $f_{-}(n) \leftrightarrow g_{-}(n)$. As to the total charge
\begin{gather}\label{eq-8.7}
	Q = \sum_{m=-\infty}^\infty [\bar{\rho}_{+}(m) + \bar{\rho}_{-}(m) + \rho_{\pm}(m)]
\end{gather}
accumulated in the whole twelve-component system (\ref{eq-6.1})--(\ref{eq-6.12}), it is seen to be conserved provided the local current $J(n)$ is the same on both spatial infinities.

The case of repulsive nonlinearity $\sigma=-1$ turns out to be more complicated, inasmuch as now the clear physical treatment of densities $\bar{\rho}_{+}(n)$ and $\bar{\rho}_{-}(n)$ as the nonpositive charge densities is enabled only under the strict limitations
\begin{gather*}%\label{eq-8.8}
	0 \leq \bar{q}_{+}(n)\bar{r}_{+}(n) < 1,
\qquad %\label{eq-8.9}
	0 \leq \bar{q}_{-}(n)\bar{r}_{-}(n) < 1 .
\end{gather*}
It is presently unknown whether or not these restrictions are globally achievable under certain type of special boundary conditions similar to those suitable for the usual semi-discrete nonlinear Schr\"{o}dinger system with the repulsive nonlinearity \cite{chubykalo-PLA-169-359, konotop-JPA-25-4037, vekslerchik-IP-8-889}. In this situation, the treatment of densities $\bar{\rho}_{+}(n)$ and $\bar{\rho}_{-}(n)$ as the nonpositive charge densities appears to be very conditional. Nevertheless, the total charge (\ref{eq-8.7})
accumulated in the whole twelve-component system \mbox{(\ref{eq-6.1})--(\ref{eq-6.12})} is obliged to be conserved even despite of its rather conditional physical treatment.

Meanwhile, the mathematical structure of expression (\ref{eq-8.2}) for the total local current $J(n)$ indicates that only four of six subsystems actually participate in the charge transportation. They are described by four pairs of fields $q_{+}(n)\leftrightarrow r_{+}(n)$, $f_{+}(n)\leftrightarrow g_{+}(n)$, $q_{-}(n)\leftrightarrow r_{-}(n)$, $f_{-}(n)\leftrightarrow g_{-}(n)$.

The six-component nonlinear integrable system is formalized by six linearly and nonlinearly coupled semi-discrete equations (\ref{eq-7.1})--(\ref{eq-7.6}) for six field functions $w_{+}(n)$, $\bar{w}_{+}(n)$, $h_{+}(n)$, $w_{-}(n)$, $\bar{w}_{-}(n)$, $h_{-}(n)$. However, in the case of six-component system, we \textit{a priori} unable to claim for the plausible pairs of presumably canonical field functions.

The fields marked by plus subscript can be treated as settled on plus labeled chain, while the fields marked by minus subscript can be treated as settled on minus labeled chain of quasi-one-dimensional regular lattice regardless of whether the system under consideration is a twelve-component or a six-component one. In this regard, the suggested twelve-component semi-discrete nonlinear integrable system (\ref{eq-6.1})--(\ref{eq-6.12}) is proved to be very prospective tool for modeling the physical properties of multicomponent essentially quasi-one-dimensional latticed objects under the action of external uniform magnetic field and external parametric drive encodeable in its coupling parameters.

\section{Conclusion}
%\label{sec9}

In our research, we proposed two novel multicomponent semi-discrete nonlinear integrable systems prospective for modeling the transport phenomena in regular quasi-one-dimensional structures of both natural and synthetic origins.

We expect the comprehensive investigation of these semi-discrete nonlinear integrable systems will be interesting both from the physical and mathematical standpoints.
Presently, the most evident open problems are (1) to construct the rigorous analytical solutions, and (2) to disclose the Hamiltonian and Poisson structures typifying the suggested semi-discrete nonlinear integrable systems.

In our opinion, the most straightforward way to obtain the explicit analytical solutions to multicomponent semi-discrete nonlinear integrable systems is based upon the Darboux--B\"{a}cklund transformation technique \cite{vakhnenko-JMP-56-033505,vakhnenko-JNMP-24-250, vakhnenko-JMP-59-053504,vakhnenko-PRE-108-024223, vakhnenko-PRSA-477-20210562, vakhnenko-EPJP-137-1176}.

The problem to establish the Hamiltonian and Poisson structures of proposed integrable systems turn out to be immensely more complicated as compared with the analogous rather nontrivial problems successfully solved in our previous works \cite{vakhnenko-JMP-56-033505, vakhnenko-JMP-57-113504,vakhnenko-JNMP-24-250, vakhnenko-JMP-59-053504, vakhnenko-LMP-108-1807}.
Some aspects of our preliminary approach to these problems are reported in Appendix~\ref{app}.

\appendix

\section{Preliminaries to Hamiltonian treatment}\label{app}

Having observed that all intersite interactions in the twelve-component semi-discrete nonlinear integrable system (\ref{eq-6.1})--(\ref{eq-6.12}) are of nearest-neighbouring
type it is reasonable to construct the system's Hamiltonian function relying upon the local conserved densities (\ref{eq-4.29})--(\ref{eq-4.32}) characterized by the same type
of couplings between the involved fields. As a result, we come to the following Hamiltonian function:
\begin{align}
 H ={}& - \alpha \sum_{m=-\infty}^\infty [q_{-}(m)\bar{r}_{+}(m)+\bar{q}_{-}(m)r_{+}(m)-\bar{q}_{-}(m)\bar{r}_{+}(m)]\nonumber\\
 & - \alpha \sum_{m=-\infty}^\infty [q_{+}(m)r_{-}(m-1)+q_{-}(m)r_{+}(m)+f_{+}(m)g_{-}(m-1)]\nonumber\\
 & - \beta \sum_{m=-\infty}^\infty [r_{-}(m)\bar{q}_{+}(m)+\bar{r}_{-}(m)q_{+}(m)-\bar{r}_{-}(m)\bar{q}_{+}(m)]\nonumber\\
 & - \beta \sum_{m=-\infty}^\infty [r_{+}(m)q_{-}(m-1)+r_{-}(m)q_{+}(m)+g_{+}(m)f_{-}(m-1)] .\label{eq-A.1}
\end{align}

According to the general rules \cite{maschke-JFI-329-923, torres-del-castillo-RMF-50-608}, the Hamiltonian dynamic equations of motion for the system under study (\ref{eq-6.1})--(\ref{eq-6.12}) must be sought in the form
\begin{gather}\label{eq-A.2}
	\frac{\mathrm{d}}{\mathrm{d}\tau}y_\lambda(n) = \sum_{\varkappa=1}^{12} \sum_{m=-\infty}^\infty J_{\lambda\varkappa}(n|m)
	\frac{\partial H}{\partial y_{\varkappa}(m)},\qquad \lambda = 1, 2, 3, \dots, 12 ,
\end{gather}
where $J_{\lambda\varkappa}(n|m)$ are the elements of skew-symmetric $J_{\varkappa\lambda}(m|n)=-J_{\lambda\varkappa}(n|m)$ symplectic matrix. These elements $J_{\lambda\varkappa}(n|m)$ are obliged to define the Poisson bracket
\begin{gather}\label{eq-A.3}
	\{F,G\} = -\sum_{n=-\infty}^\infty \sum_{\lambda=1}^{12} \sum_{\varkappa=1}^{12} \sum_{m=-\infty}^\infty \frac{\partial F}{\partial y_{\lambda}(n)} J_{\lambda\varkappa}(n|m) \frac{\partial G}{\partial y_{\varkappa}(m)}
\end{gather}
subjected to the Jacobi identity
\begin{gather}\label{eq-A.4}
	\{E,\{F,G \} \} +\{F,\{G,E \} \} +\{G,\{E,F \} \}=0 .
\end{gather}
In so doing, the set of Hamiltonian equations (\ref{eq-A.2}) acquires the form
\begin{gather}\label{eq-A.5}
	\frac{\mathrm{d}}{\mathrm{d}\tau}y_\lambda(n) = \{H, y_{\lambda}(n)\},\qquad \lambda = 1, 2, 3, \dots , 12,
\end{gather}
substantiated by the set of sixty six fundamental Poisson brackets
\begin{gather}\label{eq-A.6}
\{ y_{\lambda}(n), y_{\varkappa}(m)\} = - J_{\lambda\varkappa}(n|m) .
\end{gather}

We tried to apply the above described procedure (\ref{eq-A.2})--(\ref{eq-A.6}) to our twelve-component integrable system (\ref{eq-6.1})--(\ref{eq-6.12}) relying on the adopted Hamiltonian function (\ref{eq-A.1}) and introducing the universal notations
\begin{alignat*}{3}%\label{eq-A.7}
	&y_{1}(n) = q_{-}(n),\qquad&& y_{7}(n) = r_{-}(n),&
\\ %\label{eq-A.8}
	&y_{2}(n) = q_{+}(n),\qquad&& y_{8}(n) = r_{+}(n),&
\\ %\label{eq-A.9}
	&y_{3}(n) = \bar{q}_{-}(n),\qquad&& y_{9}(n) = \bar{r}_{-}(n),&
\\ %\label{eq-A.10}
	&y_{4}(n) = \bar{q}_{+}(n),\qquad&& y_{10}(n) = \bar{r}_{+}(n),&
\\ %\label{eq-A.11}
	&y_{5}(n) = f_{-}(n),\qquad&& y_{11}(n) = g_{-}(n),&
\\ %\label{eq-A.12}
	&y_{6}(n) = f_{+}(n),\qquad&& y_{12}(n) = g_{+}(n)&
\end{alignat*}
for the system's field functions. However, we have not managed to isolate extremally huge number of candidates on the elements $J_{\lambda\varkappa}(n|m)$ of symplectic matrix to say nothing about their verification via the Jacobi identity (\ref{eq-A.4}) or its more simple equivalents \cite {maschke-JFI-329-923, torres-del-castillo-RMF-50-608, vakhnenko-JMP-59-053504, vakhnenko-LMP-108-1807}.

The main obstacle in achieving the positive result is the generic spatial nonlocality
of inspected symplectic matrix $J_{\lambda\varkappa}(n|m)$ pronouncedly contrasting with the crucial simplification
 \begin{gather}\label{eq-A.13}
 	J_{\lambda\varkappa}(n|m) = J_{\lambda\varkappa}(n|n)\delta_{nm}
\end{gather}
typical of earlier studied systems \cite{vakhnenko-JMP-56-033505, vakhnenko-JMP-57-113504, vakhnenko-JNMP-24-250, vakhnenko-JMP-59-053504, vakhnenko-LMP-108-1807}.

\subsection*{Acknowledgements}

Oleksiy O.~Vakhnenko acknowledges support from the National Academy of Sciences of Ukraine within the Project No~0122U000887.
Vyacheslav O. Vakhnenko acknowledges support from the National Academy of Sciences of Ukraine within the Project No~0123U100182.
Oleksiy O.~Vakhnenko also acknowledges support from the Simons Foundation (USA) under the Grant SFI-PD-Ukraine-00014573.
We are grateful to referees for the constructive criticism directed to improve the quality of presented results.

\pdfbookmark[1]{References}{ref}
\begin{thebibliography}{99}
\footnotesize\itemsep=0pt

\bibitem{ablowitz-SAM-58-17}
Ablowitz M.J., Lectures on the inverse scattering transform,
 \href{https://doi.org/10.1002/sapm197858117}{\textit{Stud. Appl.
 Math.}} \textbf{58} (1978), 17--94.

\bibitem{ablowitz-PRL-31-125}
Ablowitz M.J., Kaup D.J., Newell A.C., Segur H., Nonlinear-evolution equations
 of physical significance,
 \href{https://doi.org/10.1103/PhysRevLett.31.125}{\textit{Phys. Rev. Lett.}}
 \textbf{31} (1973), 125--127.

\bibitem{ablowitz-SAM-53-249}
Ablowitz M.J., Kaup D.J., Newell A.C., Segur H., The inverse scattering
 transform-{F}ourier analysis for nonlinear problems,
 \href{https://doi.org/10.1002/sapm1974534249}{\textit{Stud. Appl. Math.}}
 \textbf{53} (1974), 249--315.

\bibitem{ablowitz-SAM-55-213}
Ablowitz M.J., Ladik J.F., A~nonlinear difference scheme and inverse
 scattering, \href{https://doi.org/10.1002/sapm1976553213}{\textit{Stud. Appl.
 Math.}} \textbf{55} (1976), 213--229.

\bibitem{ablowitz-JMP-17-1011}
Ablowitz M.J., Ladik J.F., Nonlinear differential-difference equations and
 {F}ourier analysis, \href{https://doi.org/10.1063/1.523009}{\textit{J.~Math.
 Phys.}} \textbf{17} (1976), 1011--1018.

\bibitem{ablowitz-PLA-253-287}
Ablowitz M.J., Ohta Y., Trubatch A.D., On discretizations of the vector
 nonlinear {S}chr\"odinger equation,
 \href{https://doi.org/10.1016/S0375-9601(99)00048-1}{\textit{Phys. Lett.~A}}
 \textbf{253} (1999), 287--304,
 \href{http://arxiv.org/abs/solv-int/9810014}{arXiv:solv-int/9810014}.

\bibitem{ablowitz-CSR-11-159}
Ablowitz M.J., Ohta Y., Trubatch A.D., On integrability and chaos in discrete
 systems, \href{https://doi.org/10.1016/S0960-0779(98)00280-X}{\textit{Chaos
 Solitons Fractals}} \textbf{11} (2000), 159--169,
 \href{http://arxiv.org/abs/solv-int/9810020}{arXiv:solv-int/9810020}.

\bibitem{babalic-RJP-63-114}
Babalic C.N., Integrable discretization of coupled {A}blowitz--{L}adik
 equations with branched dispersion, \textit{Romanian~J.~Phys.} \textbf{63}
 (2018), 114, 10~pages.

\bibitem{babalic-IJMPB-34-2050274}
Babalic C.N., Complete integrability and complex solitons for generalized
 {V}olterra system with branched dispersion,
 \href{https://doi.org/10.1142/S0217979220502744}{\textit{Internat.~J.~Modern
 Phys.~B}} \textbf{34} (2020), 2050274, 12~pages.

\bibitem{babalic-JPA-50-415201}
Babalic C.N., Carstea A.S., Coupled {A}blowitz--{L}adik equations with branched
 dispersion,
 \href{https://doi.org/10.1088/1751-8121/aa87a4}{\textit{J.~Phys.~A}}
 \textbf{50} (2017), 415201, 13~pages,
 \href{http://arxiv.org/abs/1705.10975}{arXiv:1705.10975}.

\bibitem{bender-RPP-70-947}
Bender C.M., Making sense of non-{H}ermitian {H}amiltonians,
 \href{https://doi.org/10.1088/0034-4885/70/6/R03}{\textit{Rep. Progr. Phys.}}
 \textbf{70} (2007), 947--1018,
 \href{http://arxiv.org/abs/hep-th/0703096}{arXiv:hep-th/0703096}.

\bibitem{carstea-JPA-48-055205}
Carstea A.S., Tokihiro T., Coupled discrete {K}d{V} equations and modular
 genetic networks,
 \href{https://doi.org/10.1088/1751-8113/48/5/055205}{\textit{J.~Phys.~A}}
 \textbf{48} (2015), 055205, 12~pages.

\bibitem{chen-AM-32-2001893}
Chen Z., Narita A., M\"ullen K., Graphene nanoribbons: On-surface synthesis and
 integration into electronic devices,
 \href{https://doi.org/10.1002/adma.202001893}{\textit{Adv. Mater.}}
 \textbf{32} (2020), 2001893, 26~pages.

\bibitem{chubykalo-PLA-169-359}
Chubykalo O.A., Konotop V.V., V\'azquez L., Vekslerchik V.E., Some features of
 the repulsive discrete nonlinear Schr\"odinger equation,
 \href{https://doi.org/10.1016/0375-9601(92)90241-D}{\textit{Phys. Let.~A}}
 \textbf{169} (1992), 359--363.

\bibitem{faddeev-HMTS-2007}
Faddeev L.D., Takhtajan L.A., Hamiltonian methods in the theory of solitons,
 \textit{Classics Math.}, \href{https://doi.org/10.1007/978-3-540-69969-9}{Springer},
 Berlin, 2007.

\bibitem{feynman-FLP3-1963}
Feynman R.P., Leighton R.B., Sands M., The {F}eynman lectures on physics.
 {V}ol.~{III}: {Q}uantum mechanics, Addison-Wesley
 Publishing Co., Reading, Mass., 1963.

\bibitem{gerdzhikov-TMP-52-676}
Gerdzhikov V.S., Ivanov M.I., {H}amiltonian structure of multicomponent
 nonlinear {S}chr\"odinger equations in difference form,
 \href{https://doi.org/10.1007/BF01027788}{\textit{Theoret. and Math. Phys}}
 \textbf{52} (1982), 676--685.

\bibitem{konotop-JPA-25-4037}
Konotop V.V., Vekslerchik V.E., On soliton creation in the nonlinear
 {S}chr\"odinger models: discrete and continuous versions,
 \href{https://doi.org/10.1088/0305-4470/25/14/018}{\textit{J.~Phys.~A}}
 \textbf{25} (1992), 4037--4042.

\bibitem{konotop-RMP-88-035002}
Konotop V.V., Yang J., Zezyulin D.A., Nonlinear waves in
 $\mathcal{PT}$-symmetric systems,
 \href{https://doi.org/10.1103/RevModPhys.88.035002}{\textit{Rev. Mod. Phys.}}
 \textbf{88} (2016), 035002, 59~pages.

\bibitem{maschke-JFI-329-923}
Maschke B.M., van~der Schaft A.J., Breedveld P.C., An~intrinsic {H}amiltonian
 formulation of network dynamics: {N}onstandard {P}oisson structures and
 gyrators,
 \href{https://doi.org/10.1016/S0016-0032(92)90049-M}{\textit{J.~Franklin
 Inst.}} \textbf{329} (1992), 923--966.

\bibitem{peierls-ZPh-80-763}
Peierls R., Zur {T}heorie des {D}iamagnetismus von {L}eitungselektronen,
 \href{https://doi.org/10.1007/BF01342591}{\textit{Z.~Phys.}} \textbf{80}
 (1933), 763--791.

\bibitem{perepelytsya-EBJ-50-759}
Perepelytsya S., Uli\v{c}n\'y J., Volkov S.N., Molecular dynamics study of the
 competitive binding of hydrogen peroxide and water molecules with {DNA}
 phosphate groups,
 \href{https://doi.org/10.1007/s00249-021-01522-8}{\textit{Eur. Biophys.~J.}}
 \textbf{50} (2021), 759--770,
 \href{http://arxiv.org/abs/2004.11882}{arXiv:2004.11882}.

\bibitem{satchanska-PM-16-1159}
Satchanska G., Davidova S., Petrov P.D., Natural and synthetic polymers for
 biomedical and environmental applications,
 \href{https://doi.org/10.3390/polym16081159}{\textit{Polymers}} \textbf{16}
 (2024), 1159, 33~pages.

\bibitem{toda-JPSJ-22-431}
Toda M., Vibration of a chain with nonlinear interaction,
 \href{https://doi.org/10.1143/JPSJ.22.431}{\textit{J.~Phys. Soc. Japan}}
 \textbf{22} (1967), 431--436.

\bibitem{toda-JPSJ-23-501}
Toda M., Wave propagation in anharmonic lattices,
 \href{https://doi.org/10.1143/JPSJ.23.501}{\textit{J.~Phys. Soc. Japan}}
 \textbf{23} (1967), 501--506.

\bibitem{toda-PhR-18-1}
Toda M., Studies of a non-linear lattice,
 \href{https://doi.org/10.1016/0370-1573(75)90018-6}{\textit{Phys. Rep.}}
 \textbf{18} (1975), 1--123.


\bibitem{torres-del-castillo-RMF-50-608}
Torres~del Castillo G.F., Vel\'azquez~Quesada M.P., Symplectic structures and
 dynamical symmetry groups, \textit{Rev. Mex. de F\'is.} \textbf{50} (2004),
 608--613.

\bibitem{tsuchida-JMP-39-4785}
Tsuchida T., Ujino H., Wadati M., Integrable semi-discretization of the coupled
 modified {K}d{V} equations,
 \href{https://doi.org/10.1063/1.532537}{\textit{J.~Math. Phys.}} \textbf{39}
 (1998), 4785--4813.

\bibitem{tsuchida-JPA-32-2239}
Tsuchida T., Ujino H., Wadati M., Integrable semi-discretization of the coupled
 nonlinear {S}chr\"odinger equations,
 \href{https://doi.org/10.1088/0305-4470/32/11/016}{\textit{J.~Phys.~A}}
 \textbf{32} (1999), 2239--2262,
 \href{http://arxiv.org/abs/solv-int/9903013}{arXiv:solv-int/9903013}.

\bibitem{tu-JPA-23-3903}
Tu G.-Z., A~trace identity and its applications to the theory of discrete
 integrable systems,
 \href{https://doi.org/10.1088/0305-4470/23/17/020}{\textit{J.~Phys.~A}}
 \textbf{23} (1990), 3903--3922.

\bibitem{vakhnenko-JPA-32-5735}
Vakhnenko O.O., Nonlinear beating excitations on ladder lattice,
 \href{https://doi.org/10.1088/0305-4470/32/30/315}{\textit{J.~Phys.~A}}
 \textbf{32} (1999), 5735--5748.

\bibitem{vakhnenko-PRE-60-R2492}
Vakhnenko O.O., Nonlinear model of intramolecular excitations on a multileg
 ladder lattice,
 \href{https://doi.org/10.1103/PhysRevE.60.R2492}{\textit{Phys. Rev.~E}}
 \textbf{60} (1999), R2492--R2495.

\bibitem{vakhnenko-PRE-77-026604}
Vakhnenko O.O., Enigma of probability amplitudes in {H}amiltonian formulation
 of integrable semidiscrete nonlinear {S}chr\"odinger systems,
 \href{https://doi.org/10.1103/PhysRevE.77.026604}{\textit{Phys. Rev.~E}}
 \textbf{77} (2008), 026604, 9~pages.

\bibitem{vakhnenko-JNMP-18-401}
Vakhnenko O.O., Semidiscrete integrable nonlinear systems generated by the new
 fourth-order spectral operator. {L}ocal conservation laws,
 \href{https://doi.org/10.1142/S1402925111001672}{\textit{J.~Nonlinear Math.
 Phys.}} \textbf{18} (2011), 401--414.

\bibitem{vakhnenko-JNMP-18-415}
Vakhnenko O.O., Semidiscrete integrable nonlinear systems generated by the new
 fourth order spectral operator. {S}ystems of obverse type,
 \href{https://doi.org/10.1142/S1402925111001684}{\textit{J.~Nonlinear Math.
 Phys.}} \textbf{18} (2011), 415--425.

\bibitem{vakhnenko-JNMP-20-606}
Vakhnenko O.O., Four-wave semidiscrete nonlinear integrable system with
 {$\mathcal{PT}$}-symmetry,
 \href{https://doi.org/10.1080/14029251.2013.865827}{\textit{J.~Nonlinear
 Math. Phys.}} \textbf{20} (2013), 606--622.

\bibitem{vakhnenko-JMP-56-033505}
Vakhnenko O.O., Nonlinear integrable model of {F}renkel-like excitations on
 a~ribbon of triangular lattice,
 \href{https://doi.org/10.1063/1.4914510}{\textit{J.~Math. Phys.}} \textbf{56}
 (2015), 033505, 21~pages.

\bibitem{vakhnenko-JMP-57-113504}
Vakhnenko O.O., Symmetry-broken canonizations of the semi-discrete integrable
 nonlinear {S}chr\"odinger system with background-controlled inter-site
 coupling, \href{https://doi.org/10.1063/1.4968244}{\textit{J.~Math. Phys.}}
 \textbf{57} (2016), 113504, 16~pages.

\bibitem{vakhnenko-JNMP-24-250}
Vakhnenko O.O., Semi-discrete integrable nonlinear {S}chr\"odinger system with
 background-controlled inter-site resonant coupling,
 \href{https://doi.org/10.1080/14029251.2017.1316011}{\textit{J.~Nonlinear
 Math. Phys.}} \textbf{24} (2017), 250--302.

\bibitem{vakhnenko-JMP-59-053504}
Vakhnenko O.O., Integrable nonlinear {S}chr\"odinger system on a lattice with
 three structural elements in the unit cell,
 \href{https://doi.org/10.1063/1.4994622}{\textit{J.~Math. Phys.}} \textbf{59}
 (2018), 053504, 25~pages.

\bibitem{vakhnenko-LMP-108-1807}
Vakhnenko O.O., Six-component semi-discrete integrable nonlinear
 {S}chr\"odinger system,
 \href{https://doi.org/10.1007/s11005-018-1049-0}{\textit{Lett. Math. Phys.}}
 \textbf{108} (2018), 1807--1824.

\bibitem{vakhnenko-PRE-108-024223}
Vakhnenko O.O., Vakhnenko V.O., Dipole-monopole criticality and chargeless half
 mode in an~integrable gauge-coupled pseudoexciton-phonon system on a~regular
 one-dimensional lattice,
 \href{https://doi.org/10.1103/physreve.108.024223}{\textit{Phys. Rev.~E}}
 \textbf{108} (2023), 024223, 8~pages.

\bibitem{vakhnenko-UJP-69-168}
Vakhnenko O.O., Vakhnenko V.O., Development and analysis of novel integrable
 nonlinear dynamical systems on quasi-one-dimensional lattices. Two-component
 nonlinear system with the on-site and spatially distributed inertial mass
 parameters,
 \href{https://doi.org/10.15407/ujpe69.3.168}{\textit{Ukr.~J.~Phys.}}
 \textbf{69} (2024), 168--178.

\bibitem{vakhnenko-UJP-69-577}
Vakhnenko O.O., Vakhnenko V.O., Development and analysis of novel integrable
 nonlinear dynamical systems on quasi-one-dimensional lattices. Parametrically
 driven nonlinear system of pseudo-excitations on a~two-leg ladder lattice,
 \href{https://doi.org/10.15407/ujpe69.8.577}{\textit{Ukr.~J.~Phys.}}
 \textbf{69} (2024), 577--590.


\bibitem{vakhnenko-PRE-61-7110}
Vakhnenko O.O., Velgakis M.J., Transverse and longitudinal dynamics of
 nonlinear intramolecular excitations on multileg ladder lattices,
 \href{https://doi.org/10.1103/PhysRevE.61.7110}{\textit{Phys. Rev.~E}}
 \textbf{61} (2000), 7110--7120.

\bibitem{vakhnenko-PRSA-477-20210562}
Vakhnenko O.O., Verchenko A.P., Nonlinear system of {$\mathcal{PT}$}-symmetric
 excitations and {T}oda vibrations integrable by the {D}arboux--{B}\"acklund
 dressing method,
 \href{https://doi.org/10.1098/rspa.2021.0562}{\textit{Proc.~R.~Soc.~A}} \textbf{477}
 (2021), 20210562, 18~pages.

\bibitem{vakhnenko-PLA-452-128460}
Vakhnenko O.O., Verchenko A.P., Branched-dispersion generalizations of
 {L}otka--{V}olterra and {A}blowitz--{L}adik nonlinear integrable systems
 revisited from the intersite coupling standpoint,
 \href{https://doi.org/10.1016/j.physleta.2022.128460}{\textit{Phys. Lett.~A}}
 \textbf{452} (2022), 128460, 6~pages.

\bibitem{vakhnenko-EPJP-137-1176}
Vakhnenko O.O., Verchenko A.P., Dipole--monopole alternative in nonlinear
 dynamics of an~integrable gauge-coupled exciton-phonon system on
 a~one-dimensional lattice,
 \href{https://doi.org/10.1140/epjp/s13360-022-03335-w}{\textit{Eur.
 Phys.~J.~Plus}} \textbf{137} (2022), 1176, 18~pages.

\bibitem{vekslerchik-IP-8-889}
Vekslerchik V.E., Konotop V.V., Discrete nonlinear {S}chr\"odinger equation
 under nonvanishing boundary conditions,
 \href{https://doi.org/10.1088/0266-5611/8/6/007}{\textit{Inverse Problems}}
 \textbf{8} (1992), 889--909.

\bibitem{wen-AML-123-107591}
Wen X.-Y., Yuan C.-L., Controllable rogue wave and mixed interaction solutions
 for the coupled {A}blowitz--{L}adik equations with branched dispersion,
 \href{https://doi.org/10.1016/j.aml.2021.107591}{\textit{Appl. Math. Lett.}}
 \textbf{123} (2022), 107591, 7~pages.

\bibitem{yu-WM-94-102500}
Yu F., Yu J., Li L., Some discrete soliton solutions and interactions for the
 coupled {A}blowitz--{L}adik equations with branched dispersion,
 \href{https://doi.org/10.1016/j.wavemoti.2019.102500}{\textit{Wave Motion}}
 \textbf{94} (2020), 102500, 11~pages.

\bibitem{yuan-CSF-168-113180}
Yuan C., Yang H., Meng X., Tian Y., Zhou Q., Liu W., Modulational instability
 and discrete rogue waves with adjustable positions for a two-component
 higher-order {A}blowitz--{L}adik system associated with {$4\times 4$} {L}ax
 pair, \href{https://doi.org/10.1016/j.chaos.2023.113180}{\textit{Chaos
 Solitons Fractals}} \textbf{168} (2023), 113180, 12~pages.

\end{thebibliography}\LastPageEnding

\end{document} 